\documentclass[12pt]{amsart}
\usepackage[utf8]{inputenc}
\usepackage{amssymb}
\usepackage{amsmath}
\usepackage{amsthm}
\usepackage{graphicx}
\usepackage{amscd}
\usepackage{epic}
%\usepackage{comment}
%\usepackage[alphabetic]{amsrefs}
%\usepackage{setspace}
\usepackage{xypic}
\usepackage[colorlinks, citecolor=blue]{hyperref}
\usepackage[left=1.5in,right=1.5in,top=1.5in,bottom=1.5in]{geometry}
\usepackage{soul}
%\usepackage{tikz}
%\usepackage[pagewise]{lineno}\linenumbers
%\usepackage[inline]{showlabels}
\usepackage{comment}
\usepackage{enumitem}
\usepackage{mathrsfs}

%\input{diagrams.sty}
%\topmargin=-5mm
%\topmargin=0mm
%\textwidth=150mm \textheight=220mm \oddsidemargin=5mm
%%%%%%%%%%%%%%%%%%%%%%%%%%%%%%%%%%%%%%%%%%%%%%%%%%%%%%%%%%%%%

\theoremstyle{plain}
\newtheorem{theorem}{Theorem}[section]
\newtheorem{corollary}[theorem]{Corollary}
\newtheorem{lemma}[theorem]{Lemma}
\newtheorem{conjecture}[theorem]{Conjecture}

\newtheorem{example}[theorem]{Example}
\newtheorem{exc}[theorem]{Exercise}
\newtheorem{proposition}[theorem]{Proposition}
\newtheorem{definition-lemma}[theorem]{Definition-Lemma}
\theoremstyle{remark}
\newtheorem{remark}[theorem]{Remark}

\theoremstyle{definition}
\newtheorem{question}[theorem]{Question}
\newtheorem{convention}[theorem]{Convention}

\newtheorem{axiom}[theorem]{Axiom}
%%%  JK commands
\newcommand{\Lk}[0]{\operatorname{Lk}}
\newcommand{\lk}[0]{\operatorname{Lk}}
%\newcommand{\st}[0]{\operatorname{St}}
\newcommand{\cst}[0]{\overline{\operatorname{St}}}
\newcommand{\St}[0]{\overline{\operatorname{St}}}
\newcommand{\ex}[0]{\operatorname{Ex}}
\newcommand{\spec}[0]{\operatorname{Spec}}
\newcommand{\red}[0]{\operatorname{red}}
\newcommand{\supp}[0]{\operatorname{Supp}}
\newcommand{\map}{\dasharrow}
\newcommand{\cone}[0]{\operatorname{Cone}}
\newcommand{\dlt}[0]{\operatorname{dlt}}
\newcommand{\Aut}[0]{\operatorname{Aut}}
\newcommand{\Spr}[0]{\operatorname{Spr}}
\newcommand{\Tr}[0]{\operatorname{Tr}}
\newcommand{\Pic}[0]{\operatorname{Pic}}

%%% end JK commands
\newcommand{\wh}{\widehat}
\newcommand{\wt}{\widetilde}
\def\ep{\varepsilon}
\def\ddbar{\partial\bar\partial}
\def\ve{\varepsilon}
\def\mcS{\mathcal{S}}
\def\NE{\overline{\operatorname{NE}}}
\def\NS{\operatorname{NS}}
\def\NF{\operatorname{\overline{NF}}}
\def\NM{\operatorname{\overline{NM}}}
\def\BC{\operatorname{BC}}
\def\NA{\overline{\operatorname{NA}}}
\def\Null{\operatorname{Null}}
\newcommand{\<}{\leq}
\def\>{\geq}
\newcommand{\mbQ}{\mathbb{Q}}
\newcommand{\mbR}{\mathbb{R}}
\newcommand{\CC}{\mathbb{C}}
\newcommand{\cC}{{\mathcal{C}}}

\def\mcO{\mathcal{O}}
\newcommand{\num}{\equiv}
\newcommand{\Fut}{{\rm DF}}
\newcommand{\MO}{\mathcal{O}}
\newcommand{\dbr}{\bar{\partial}}
\newcommand{\bX}{{\mathbbl{X} }}
\newcommand{\Z}{{\mathbb{Z} }}
\newcommand{\bL}{{\mathbb{L}}}
\newcommand{\sM}{{\mathcal{M}}}
\newcommand{\sL}{{\mathcal{L}}}
\newcommand{\mJ}{{\mathcal{J}}}
\newcommand{\mX}{{\mathcal{X}}}
\newcommand{\mY}{{\mathcal{Y}}}
\newcommand{\mZ}{{\mathcal{Z}}}
\newcommand{\kX}{{\mathfrak{X}}}
\newcommand{\kY}{{\mathfrak{Y}}}
\newcommand{\bM}{{\bar{\mathcal{M}}}}
\newcommand{\bH}{{\bar{\mathcal{H}}}}
\newcommand{\bY}{{\bar{\mathcal{Y}}}}
\newcommand{\DMR}{{\mathcal{DMR}}}
\newcommand{\DR}{{\mathcal{DR}}}
%\newcommand{\Star}{{\rm St}}
\newcommand{\D}{{\mathcal{D}}}
\newcommand{\OO}{{\mathcal{O}}}
\newcommand{\Q}{{\mathbb{Q}}}
\newcommand{\C}{{\mathbb{C}}}
\newcommand{\R}{{\mathbb{R}}}
\newcommand{\mop}{{\mathcal{O}_{\mathbb{P}^1}(1)}}
\newcommand{\mot}{{\mathcal{O}_{\mathbb{P}^1}(-1)}}
%\newcommand{\Lk}{{\rm Lk}}
%\newcommand{\St}{{\rm St}}
%\newcommand{\st}{{\rm st}}
\newcommand{\mld}{{\rm mld}}
\newcommand{\mult}{{\rm mult}}
\newcommand{\Supp}{{\rm Supp}}
\newcommand{\Hom}{{\rm Hom}}
\newcommand{\sm}{{\rm sm}}
\newcommand{\stringy}{{\rm st}}
\newcommand{\topo}{{\rm top}}
\newcommand{\ord}{{\rm ord}}
\newcommand{\moti}{{\rm mot}}
\newcommand{\Link}{{\rm Link}}





%\newcommand{\mbR}{\mathbb{R}}
\newcommand{\mbC}{\mathbb{C}}
\newcommand{\mbZ}{\mathbb{Z}}
%\newcommand{\mbQ}{\mathbb{Q}}
\newcommand{\mbK}{\mathbb{K}}
\newcommand{\bir}{\dashrightarrow}
\newcommand{\vphi}{\varphi}
\def\injective{\hookrightarrow}
\def\surjective{\twoheadrightarrow}
\newcommand{\del}{\partial}
\def\lrd{\lfloor}
\def\rrd{\rfloor}


\def\mbF{\mathbb{F}}
\def\mbN{\mathbb{N}}
\def\mbP{\mathbb{P}}
\def\>{\geq}
\def\ve{\varepsilon}
\def\mcA{\mathcal{A}}
\def\mcO{\mathcal{O}}
\def\mcB{\mathcal{B}}
\def\mcC{\mathcal{C}}
\def\mcD{\mathcal{D}}
\def\mcE{\mathcal{E}}
\def\mcF{\mathcal{F}}
\def\mcG{\mathcal{G}}
\def\mcH{\mathcal{H}}
\def\mcK{\mathcal{K}}
\def\mcL{\mathcal{L}}
\def\mcM{\mathcal{M}}
\def\mcN{\mathcal{N}}
\def\mcP{\mathcal{P}}
\def\mcR{\mathcal{R}}
\def\mcS{\mathcal{S}}
\def\mcT{\mathcal{T}}
\def\mcI{\mathcal{I}}
\def\mcJ{\mathcal{J}}
\def\mcV{\mathcal{V}}
\def\mcW{\mathcal{W}}
\def\mcZ{\mathcal{Z}}
\def\msE{\mathscr{E}}
\def\msF{\mathscr{F}}
\def\msG{\mathscr{G}}
\def\msH{\mathscr{H}}
\def\msI{\mathscr{I}}
\def\msL{\mathscr{L}}
\def\msM{\mathscr{M}}
\def\msN{\mathscr{N}}

\def\eps{\epsilon}
\def\bfP{\textbf{P}}
\def\bfQ{\textbf{Q}}
\def\bfR{\mathbf{R}}

\def\lru{\lceil}
\def\rru{\rceil}
\def\lrd{\lfloor}
\def\rrd{\rfloor}


\def\bbeta{\boldsymbol{\beta}}
\def\bomega{\boldsymbol{\omega}}
%\newcommand{\mot}{{\rm mot}}
%\begin{CJK*}{GBK}{Song}
% \theoremstyle{plain}
% \newtheorem{thm}[thm]{Theorem}%[section]
% \newtheorem{prop}[thm]{Proposition}%[section]
% \newtheorem{cor}[thm]{Corollary}
% \newtheorem{lem}[thm]{Lemma}
% \newtheorem{conj}[thm]{Conjecture}
% \theoremstyle{definition}
% \newtheorem{defn}{Definition}%[section]
% \newtheorem{rem}{Remark}
% \newtheorem{que}{Question}
% \newtheorem{exmp}{Example}%[section]
% \newtheorem{claim}{Claim}
% \newtheorem{assum}{Assumption}



\def\pt{\operatorname{pt}}
\def\Ex{\operatorname{Ex}}
\def\dim{\operatorname{dim}}
\def\codim{\operatorname{codim}}
\def\sing{\operatorname{\textsubscript{sing}}}
\def\sm{\operatorname{\textsubscript{sm}}}
\def\NS{\operatorname{NS}}
\def\NA{\operatorname{\overline{NA}}}
\def\NAA{\operatorname{NA}}
\def\NE{\operatorname{\overline{NE}}}
\def\Nef{\operatorname{Nef}}
\def\Im{\operatorname{Im}}
\def\Chow{\operatorname{Chow}}
\def\EnK{\operatorname{E_{nK}}}
\def\rank{\operatorname{rank}}
\def\Div{\operatorname{Div}}
\def\lct{\operatorname{lct}}
\def\LCT{\operatorname{LCT}}
\def\Bs{\operatorname{Bs}}
\def\Zero{\operatorname{Zero}}
\def\Alb{\operatorname{Alb}}
\def\Pic{\operatorname{Pic}}
\def\Proj{\operatorname{Proj}}
\def\Spec{\operatorname{Spec}}

\def\ddbar{\partial\bar{\partial}}
%\newtheorem{thm}{Theorem}[section]
%\newtheorem{mainthm}[thm]{Main Theorem}
%\newtheorem{question}[thm]{Question}
%\newtheorem{lem}[thm]{Lemma}
%\newtheorem{cor}[thm]{Corollary}
%\newtheorem{cors}[thm]{Corollaries}
%\newtheorem{prop}[thm]{Proposition}
%\newtheorem{crit}[thm]{Criterion}
%\newtheorem{conj}[thm]{Conjecture}
%\newtheorem{principle}[thm]{Principle} %!!!!!!!!!!!!!!!!!!!!!!
%\newtheorem{complem}[thm]{Complement}%!!!!!!!!!!!!!!!!!!!!!!
%\newtheorem{pitfall}[thm]{Pitfall}%!!!!!!!!!!!!!!!!!!!!!!
\theoremstyle{definition}
\newtheorem{definition}[theorem]{Definition}
\theoremstyle{definition}
\newtheorem{notation}[theorem]{Notation}
\numberwithin{equation}{section}

% \theoremstyle{definition}
% \newtheorem{claim}[theorem]{Claim}



%\theoremstyle{definition}
%\newtheorem{defn}[thm]{Definition}
%\newtheorem{condition}[thm]{Condition}
%\newtheorem{say}[thm]{}
%\newtheorem{exmp}[thm]{Example}
%\newtheorem{hint}[thm]{Hint}
%\newtheorem{exrc}[thm]{Exercise}
%\newtheorem{prob}[thm]{Problem}
%\newtheorem{const}[thm]{Construction}   %!!!!!!!!!!!!!!!!
%\newtheorem{ques}[thm]{Question}    %!!!!!!!!!!!!!!!!!!!!
%\newtheorem{que}[thm]{Question}    %!!!!!!!!!!!!!!!!!!!!
%\newtheorem{alg}[thm]{Algorithm}
%\newtheorem{rem}[thm]{Remark}
%\newtheorem{aside}[thm]{Aside}
%\renewcommand{\theremark}{}
%\newtheorem{note}[thm]{Note}            %\renewcommand{\thenote}{}
%\newtheorem{summ}[thm]{Summary}         %\renewcommand{\thesumm}{}
%\newtheorem*{ack}{Acknowledgments}      % \renewcommand{\theack}{}
%\newtheorem{notation}[thm]{Notation}
%\newtheorem{warning}[thm]{Warning}
%\newtheorem{defn-thm}[thm]{Definition--Theorem}  %!!!!!!!!!!!!!!!!!!!!!!!!
%\newtheorem{defn-prop}[thm]{Definition--Proposition}  %!!!!!!!!!!!!!!!!!!!!!!!!
%\newtheorem{defn-lem}[thm]{Definition--Lemma}  %!!!!!!!!!!!!!!!!!!!!!!!!
%\newtheorem{convention}[thm]{Convention}  %!!!!!!!!!!!!!!!!!!!!!!!!!!!
%\newtheorem{keyidea}[thm]{Key idea}
%\newtheorem{expl}[thm]{Explanation}
%\newtheorem{meth}[thm]{Method}
%\newtheorem{assumption}[thm]{Assumption}
%\newtheorem{assertion}[thm]{Assertion}
%\newtheorem{comp}[thm]{Complement}

\theoremstyle{remark}
%\newtheorem{theorem}[Theorem]
\newtheorem{claim}[theorem]{Claim}
\newtheorem{case}{Case}
\newtheorem{subcase}{Subcase}
\newtheorem{step}{Step}
\newtheorem{approach}{Approach}
%\newtheorem{principle}{Principle}
\newtheorem{fact}{Fact}
\newtheorem{subsay}{}

\title{On the Minimal Model Program for K\"ahler 3-folds}
%[On Extremal Contractions for the K\"ahler MMP]{On Extremal Contractions for the K\"ahler MMP}
%{On the existence of flipping and divisorial contractions}


\author{Omprokash Das}
\address{School of Mathematics\\
Tata Institute of Fundamental Research\\
Homi Bhabha Road, Navy Nagar\\
Colaba, Mumbai 400005}
\email{omdas@math.tifr.res.in}
\email{omprokash@gmail.com}
\thanks{Omprokash Das was partially supported by the Start--Up Research Grant(SRG), Grant No. \# SRG/2020/000348 of the Science and Engineering Research Board (SERB), Govt. Of India.}

\author{Christopher Hacon}
\address{Department of Mathematics\\
University of Utah\\
155 S 1400 E\\
Salt Lake City, Utah 84112, USA}
\email{hacon@math.utah.edu}
\thanks{Christopher Hacon was partially supported by the NSF research grants no: DMS-
1952522, DMS-1801851, DMS-58503413.}

\date{}

\begin{document}
\maketitle

\begin{abstract}
    In this article we prove the existence of pl-flipping and divisorial contractions and pl flips in dimension $n$ for compact K\"ahler varieties, assuming results of the minimal model program in dimension $n-1$. We also give a self contained proof of the cone theorem, the existence of flipping and divisorial contractions, of flips and minimal models in dimension 3.
\end{abstract}

\tableofcontents

\section{Introduction}
In recent years there has been substantial progress towards the minimal model program for K\"ahler varieties.
The 3-fold minimal model program is established in increasing generality in \cite{CP97}, \cite{HP15}, \cite{Pet98}, \cite{Pet01}, \cite{HP16}, \cite{CHP16}, \cite{DH20}, \cite{DO23}, and \cite{DHY23}.
Further progress was then made in \cite{Fuj22}, \cite{DHP22}, and \cite{LM22}, where the minimal model program for projective morphisms is established in arbitrary dimensions. Note that flipping contractions are, by definition, projective morphisms and hence the above references imply the existence of flips in all dimensions. Unluckily, we do not have a K\"ahler analog of the base point free and cone theorems and hence this is the main difficulty in proving the minimal model program in higher dimensions. 

In dimension 3, the existence of flipping and divisorial contractions for terminal pairs was claimed in \cite{HP16}. Flipping contractions and divisorial contractions to a point for klt 3-folds follow from \cite{CHP16} and klt divisorial contractions to a curve follow from \cite{DH20}.
Unluckily, the proof of the existence of divisorial contractions to a curve  claimed in \cite{HP16} seems to be incomplete (and relies on \cite{AT84} which also seems to contain an error). To make matters more complicated, the arguments of \cite{CHP16} heavily rely on those of \cite{HP16}, and \cite{DH20} also indirectly relies on  \cite{HP16} and \cite{CHP16}.

 The purpose of this paper is to give a brief, unified, and mostly self contained treatment of the minimal model program for klt K\"ahler 3-folds. In particular, we give new proofs of the cone theorem, and the existence of flipping contractions and divisorial contractions for K\"ahler klt 3-folds.  In the process, we also establish an inductive framework that we expect to be useful in proving the existence of flipping contractions and divisorial contractions for higher dimensional K\"ahler klt pairs. The strategy of this paper is primarily inspired by the approaches of \cite{BCHM10}, \cite{DH20}, and \cite{DHP22}. %Another advantage of our approach is to give a new proof of the existence of 3-fold klt flipping contractions, which is independent (and we believe simplifies) the proof given in \cite{CHP16}. 

We will now describe the main results contained in this paper. We begin by proving the existence of pl-divisorial contractions and pl-flipping contractions and the corresponding flips in dimension $n$ assuming a base point free result in dimension $n-1$.
\begin{theorem}\label{t-pl}
    Let $(X,S+B)$ be a $\Q$-factorial compact K\"ahler plt pair of arbitrary dimension, $\pi :S\to T$ a contraction such that $-(K_X+S+B)|_S$ and $-S|_S$ are $\pi$-ample. Then there exists a projective bimeromorphic morphism $p : X \to Z$ with connected fibers such that $p|_S = \pi$
and $p|_{X\setminus S}$ is an isomorphism. If $p$ is of flipping type, then the flip $p^+:X^+\to Z$ exists.
\end{theorem}
In fact we prove a more general version of this results that applies to generalized plt pairs, see Theorems \ref{t-ext} and \ref{t-flip}.
Note that the existence of flips for plt pairs was shown in \cite{Fuj22} and \cite{DHP22}.
In order to apply the above result, we must guarantee the existence of the morphism $\pi :S\to T$. If there exists a K\"ahler form $\omega$ such that $\alpha =[K_X+S+B+\omega]$ is nef and $R=\alpha ^\perp\cap \overline{\rm NA}(X)$ is an extremal ray, then we must show that $R$ can be contracted in $S$.
To this end, we consider the pair $K_S+B_S+\omega _S=(K_S+S+B+\omega)|_S$ induced by adjunction, and we aim to prove that there is a morphism $\pi:S\to T$ such that $K_S+B_S+\omega _S\equiv \pi ^* \omega _T$ where $\omega _T$ is a K\"ahler form on $T$. 
In dimension $\dim S \leq 3$, $\pi$ exists by \cite{DHY23} and therefore Theorem \ref{t-pl} holds in dimension $\dim X\leq 4$. (We remark that the results of \cite{DHY23} only rely on the 3-fold MMP established in this paper.)
We also note that if such a morphism $\pi$ exists, then $-(K_S+B_S)\equiv _T \omega_S$ is ample over $T$ and hence this morphism is projective and so, a posteriori, the existence of $\pi$ is implied by the usual base point free theorem.

We now turn our attention to the full minimal model program for klt K\"ahler 3-folds. We first show the existence of flipping and divisorial contractions (see Theorems \ref{t-flipcont} and \ref{t-divcont}).
\begin{theorem}\label{thm:contractions}
    Let $(X,B)$ be a $\mbQ$-factorial compact K\"ahler 3-fold klt pair and $\beta$ a K\"ahler form such that $\alpha=[K_X+B+\beta]$ is nef and big and $\alpha ^\perp \cap \overline {\rm NA}(X)=R$ is an extremal ray. Then the there exists a projective bimeromorphic morphism $f:X\to Z$ with connected fibers such that 
    $\alpha \equiv f^*\alpha _Z$, where $\alpha _Z$ is a K\"ahler class on $Z$. Moreover, if $f$ is a divisorial contraction, then $(Z, f_*B)$ has $\mbQ$-factorial klt singularities. 
\end{theorem}
Using this result, we then prove the cone theorem and the existence of minimal models when $K_X+B$ is pseudo-effective. 

\begin{theorem}  \label{t-cone}  Let $(X,B)$ be a $\mbQ$-factorial compact K\"ahler $3$-fold klt pair such that $K_X+B$ is pseudo-effective. Then there are at most countably many rational curves $\{\Gamma _i\}_{i\in I}$ such that $-(K_X+B)\cdot \Gamma _i\leq 6$ for all $i\in I$ and 
 \[\overline{\rm NA}(X)=\overline{\rm NA}(X)_{(K_X+B) \geq 0}+\sum _{i\in I}\mathbb R ^+[\Gamma _i].\]
 % Moreover, if $K_X+B$ is big, then $I$ is finite.
 \end{theorem}
\begin{theorem}\label{t-mmp-mfs}
    Let $(X,B)$ be a $\Q$-factorial compact K\"ahler $3$-fold klt pair such that $K_X+B$ is pseudo-effective. Then there exists a minimal model program \[X=X_0\dasharrow X_1\dasharrow X_2\dasharrow \ldots \dasharrow X_m\] where $K_{X_m}+B_m$ is nef.
  
\end{theorem}
The existence of Mori-fiber spaces also follows, but we do not pursue it here as it also relies on the results of \cite{Bru06}, \cite{HP15}, and \cite{DH20}.\\


\begin{remark}
Theorems 10.14, 10.16 and 10.17 in \cite{DO23} also rely on the contraction results of \cite{HP16} and \cite{CHP16}, however these theorems are also implied by Theorem \ref{thm:contractions} (more specifically see Corollary \ref{cor:dlt-contraction}), and so the results of \cite{DO23} are unaffected.


\end{remark}



The article is organized in the following manner. In Section \ref{sec:prelim} we recall some definitions and provide appropriate references, in Section \ref{sec:point-contractions} prove some contraction theorems when the exceptional locus maps to finitely many points. Section \ref{sec:main-contractions} is the core part of this article, in this section we prove Theorems \ref{t-pl} and \ref{thm:contractions}. Finally, in Section \ref{sec:cone-and-mm} we prove the cone Theorem \ref{t-cone} and the existence of minimal model, Theorem \ref{t-mmp-mfs}.\\

\noindent
{\bf Acknowledgement } We would like to thank Paolo Cascini, Mihai P\u{a}un, J\'anos Koll\'ar, and Matei Toma for useful comments, references and suggestions. % for  




\section{Preliminaries}\label{sec:prelim}
We will follow the usual conventions of the minimal model program. In particular, we refer the reader to \cite[Chapter 2]{KM98} for the definition of pairs and their singularities (klt, lc, plt etc.), \cite[Subsection 2.1]{DHY23} for the definitions of generalized pairs and relates singularities (glc, gklt, gdlt, etc.), \cite[Chapters 2, 3, 11]{Fuj22} for many key concepts of the minimal model program for analytic varieties, to \cite[Definition 2.8]{DHP22} for the definitions of nef, minimal and log terminal models, \cite{HP16} for a discussion of Bott-Chern cohomology $H^{1,1}_{\rm BC}(X)$, the  K\"ahler and nef cones $\mathcal K\subset \bar {\mathcal K}$, the Mori cone and the cone of positive closed currents $\overline{\rm NE}(X)\subset \overline {\rm NA}(X)\subset N_1(X)$.  Analytic pl-flipping contractions and pl-flips are defined in \cite[Chapter 15]{Fuj22}.



\begin{definition}\label{def:mod-nef}
Let $X$ be a normal $\mbQ$-factorial compact analytic variety and $\alpha\in H^{1,1}_{\BC}(X)$ is a pseudo-effective class. Then the \emph{negative part} $N(\alpha)$ of the Boucksom-Zariski decomposition of $\alpha$ is an effective $\mbR$-divisor (see \cite[Defintion A.6]{DHY23} and \cite[Definition 3.7]{Bou04}). We note that in \cite[Definition  A.6]{DHY23} it is assumed that $X$ is K\"ahler, however, it was never used. Since $X$ is $\mbQ$-factorial, $N(\alpha)$ is $\mbR$-Cartier, and so $\alpha-[N(\alpha)]$ is a $(1,1)$ class in $H^{1,1}_{\BC}(X)$. We define $P(\alpha):=\alpha-[N(\alpha)]$ and call it the \emph{positive part} of $\alpha$. Lemma \ref{lem:positive-part} below shows that $P(\alpha)$ is a modified nef class (see \cite[Definition 2.2(ii)]{Bou04}) when $X$ is in Fujiki's class $\mcC$. We call $\alpha=[N(\alpha)]+P(\alpha)$, the \emph{ Boucksom-Zariski decomposition} of $\alpha$.  
\end{definition}

\begin{lemma}\label{lem:positive-part}
    Let $X$ be a normal $\mbQ$-factorial compact analytic variety in Fujiki's class $\mcC$ and $\alpha\in H^{1,1}_{\BC}(X)$ a pseudo-effective class. Then the positive part $P(\alpha)$ of the Boucksom-Zariski decomposition of $\alpha$ defined above is a modified nef class.
\end{lemma}

\begin{proof}
    Let $f:Y\to X$ be a resolution of singularities of $X$ such that $Y$ is a K\"ahler manifold and $f^*\alpha=N(f^*\alpha)+P(f^*\alpha)$ the Boucksom-Zariski decomposition of $f^*\alpha$ as in \cite[Defintion 3.7]{Bou04}. Then from \cite[Definition A.6]{DHY23} it follows that $P(\alpha)=f_*(P(f^*\alpha))$. Note that $P(f^*\alpha)$ is modified nef by \cite[Proposition 3.8(i)]{Bou04}. Let $\beta$ be a modified K\"ahler $(1,1)$ class on $X$; such a class exists since $Y$ is K\"ahler. It is suffices to show that $P(\alpha)+\beta$ is modified K\"ahler. By \cite[Proposition 2.3]{Bou04} there is a K\"ahler class $\omega_Y$ on $Y$ such that $f_*\omega_Y=\beta$. Then $P(\alpha)+\beta=f_*(P(f^*\alpha)+\omega_Y)$. Since the modified nef cone is the closure of the modified K\"ahler cone (see \cite[\S 2.7]{Bou04}), it follows that $P(f^*\alpha)+\omega_Y$ is modified K\"ahler on $Y$, and hence $P(\alpha)+\beta$ is modified K\"ahler by \cite[Proposition 2.3]{Bou04}.  
\end{proof}~\\

\begin{definition}
Let $X$ be a normal analytic variety and $T$ a closed positive $(1, 1)$ current on $X$ with local potentials. Then for each $c>0$ we define
\[ 
E_c:=\{x\in X\;:\; \nu(T, x)>c\},
\]
where $\nu(T, x)$ denote the Lelong number of $T$ at $x\in X$.\\ 
This is an analytic subset of $X$ by \cite{Siu74}. We define 
\[
E_+(T)=\cup_{c\in \mbR^+}E_c.
\]
A closed positive $(1,1)$ current $T$ on $X$ with local potentials is called a K\"ahler current if $T\>\omega$ for some Hermitian form $\omega$ on $X$. A $(1,1)$ class $\alpha$ is called big if it contains a K\"ahler current. If $\alpha\in H^{1,1}_{\BC}(X)$ is a nef and big class, then we define the \emph{non-K\"ahler} locus of $\alpha$ as follows:
\[
E_{nK}(\alpha ):=\cap_{T\in \alpha} E_+(T),
\]
where $T$ runs over all K\"ahler currents contained in the class $\alpha$.
\end{definition}

\begin{lemma}\label{lem:non-kahler-locus}
Let $f:Y\to X$ be a proper bimeromorphic morphism of normal compact analytic varieties in Fujiuki's class $\mcC$ and with rational singularities. Let $\alpha$ be a $(1,1)$ big class on $X$. Then $E_{nK}(f^*\alpha)=f^{-1}(E_{nK}(\alpha))\cup\Ex(f)$.
\end{lemma}

\begin{proof}
First we will show that $f^{-1}(E_{nK}(\alpha))\cup\Ex(f)\subset E_{nK}(f^*\alpha)$. So,  supposing  that $y\notin E_{nK}(f^*\alpha)$, we must show that
$y\notin f^{-1}(E_{nK}(\alpha))\cup\Ex(f)$. Since $y\notin E_{nK}(f^*\alpha)$, then there is a K\"ahler current $T\in f^*\alpha$ such that $T$ is a smooth form near $y$. 
Suppose that  $y\in \Ex(f)$. Write $T=f^*\theta+i\ddbar\vphi$, where $\theta$ is a smooth $(1,1)$ form on $X$ with local potentials representing the class $\alpha$ and $\vphi$ is a quasi-psh function on $Y$. Since $T$ is a K\"ahler current, there is a Hermitian form $\omega_Y$ such that $\theta+i\ddbar\vphi\>\omega_Y$. Let $E$ be an irreducible component of $f^{-1}(f(y))$ containing $y$ and $\widetilde E\to E$ is the normalization. Since $T$ is smooth near $y$, $E$ is not contained in the locus $\{y\in Y\;|\; \vphi(y)=-\infty\}$, in particular, $T|_{\widetilde E}$ is well defined and we have 
\begin{equation}
    T|_{\widetilde E}=(f^*\theta+i\ddbar\vphi)|_{\widetilde E}=i\ddbar(\vphi|_{\widetilde E})\>\omega_Y|_{\widetilde E}.
\end{equation}
Thus $\vphi|_{\widetilde E}$ is a strictly psh function. Since $\widetilde E$ is compact, from the maximal principle of psh functions it follows that $\vphi|_{\widetilde E}$ is a constant function, which is a contradiction to the above inequality. In particular, $y\not\in\Ex(f)$, and so
$f$ is an isomorphism near $y\in Y$. By \cite[Lemma 3.4]{HP16}, $S:=f_*T$ is a closed positive $(1,1)$ current on $X$ with local potentials such that $[S]=\alpha$; clearly $S$ is a K\"ahler current. Since $f$ is isomorphism near $y\in Y$, it follows that $S$ is a smooth form near $x=f(y)\in X$. 
In particular, $x\not\in E_{nK}(\alpha)$, and the proof of the inclusion  $f^{-1}(E_{nK}(\alpha))\cup\Ex(f)\subset E_{nK}(f^*\alpha)$ is complete.\\ 

For the reverse inequality, suppose that $y\not\in f^{-1}(E_{nK}(\alpha))\cup\Ex(f)$. Then $f$ is an isomorphism near $y$ and $x=f(y)\not\in E_{nK}(\alpha)$. Thus there is a K\"ahler current $T\in \alpha$ such that $T$ is smooth near $x$, in particular, $f^*T\in f^*\alpha$ is K\"ahler current on $Y$ which is smooth near $y$. Hence $y\not\in E_{nK}(f^*\alpha)$. This completes our proof. 


\end{proof}


\begin{lemma}\label{lem:null-locus}
    Let $f:Y\to X$ be a proper bimeromorphic morphism between two normal compact analytic varieties and $\alpha$ is an $(1,1)$ nef and big class on $X$. Then $\Null(f^*\alpha)=f^{-1}(\Null(\alpha))\cup \Ex(f)$. 
\end{lemma}

\begin{proof}
First we will show that $\Null(f^*\alpha)\subset f^{-1}(\Null(\alpha))\cup \Ex(f)$. To that end assume that $y\notin f^{-1}(\Null(\alpha))\cup \Ex(f)$. Then $f$ is an isomorphism near $y\in Y$. 
Assume by contradiction that $y\in \Null(f^*\alpha)$, then there is a subvariety $V\subset Y$ containing $y$ such that $(f^*\alpha)^{\dim V}\cdot V=0$. Then by the projection formula we have $\alpha^{\dim V}\cdot f_*V=0$. However, since $f$ is an ismorphism near $y$, $f|_V:V:\to f(V)$ is bimeromorphic, and in particular, $f_*V\neq 0$. Therefore $f(y)\in f(V)\subset \Null(\alpha)$, and so $y\in f^{-1}(\Null(\alpha))$, which is a contradiction.\\ 
Now we will show the reverse inclusion. Let $y\in \Ex(f)$ and $E$ an irreducible component of $f^{-1}(f(y))$ containing $y$. Then clearly, $(f^*\alpha)^{\dim E}\cdot E=0$, and thus $y\in \Null(f^*\alpha)$. Now choose $y\in f^{-1}(\Null(\alpha))\setminus\Ex(f)$; then $f$ is an isomorphism near $y$. Since $f(y)\in \Null(\alpha)\setminus f(\Ex(f))$, there is a subvariety $f(y)\in W\subset X$ not contained in $f(\Ex(f))$ such that $\alpha^{\dim W}\cdot W=0$. Let $V\subset Y$ be the strict transform of $W$ under $f$. Then by the projection formula we have $(f^*\alpha)^{\dim V}\cdot V=0$; since $y\in V$, it follows that $y\in \Null(f^*\alpha)$. This completes the proof.
         
\end{proof}




\begin{definition}\label{def:modified-kahler}
    Let $X$ be a normal compact analytic variety and $\alpha\in H^{1,1}_{\BC}(X)$. We say that $\alpha$ is a \emph{modified K\"ahler class} if there is a closed positive $(1,1)$ current $T$ with local potentials in $\alpha$ such that the Lelong numbers satisfy $\nu(T, P)=0$ for all prime Weil divisors $P\subset X$.  
\end{definition}

We prove the following singular and weaker version of \cite[Proposition 2.3]{Bou04}
\begin{lemma}\label{lem:modified-kahlerle}
Let $X$ be a normal compact analytic variety in Fujiki's class $\mcC$ and $\alpha\in H^{1,1}_{\BC}(X)$. Then the following hold:
\begin{enumerate}
    \item If there is a proper bimeromorphic morphism $f:Y\to X$ from a K\"ahler manifold $Y$ and a K\"ahler class $\omega_Y$ such that $f_*\omega_Y=\alpha$, then $\alpha$ is a modified K\"ahler class.
    \item If $\alpha$ modified K\"ahler and nef, then there is a proper bimeromorphic morphism $f:Y\to X$ from a K\"ahler manifold $Y$ and a K\"ahler class $\omega_Y$ such that $f_*\omega_Y=\alpha$.
\end{enumerate}
\end{lemma}

\begin{proof}
    (1) follows from the same argument as \cite[Proposition 2.3]{Bou04}.\\
    
    (2) First note that, since $\alpha$ is nef and big and $X$ is in Fujiki's class $\mcC$, $\alpha^{\dim X}>0$. Thus $\Null(\alpha)\neq X$. Moreover, since $\alpha$ is modified K\"ahler, the non-K\"ahler locus $E_{nK}(\alpha)$ doesn't contain any divisor. Let $f:Y\to X$ be a resolution of singularities of $X$. Then by \cite[Theorem 1.1]{CT15}, $E_{nK}(f^*\alpha)=\Null(f^*\alpha)$, and so from Lemmas \ref{lem:non-kahler-locus} and \ref{lem:null-locus} it follows that $E_{nK}$ and $\Null(\alpha)$ are both contained in some analytic subset of $X$ of codimension $\>2$. Let $Z\subset X$ be such an analytic set, i.e. $E_{nK}\cup\Null(\alpha)\subset Z$ and $\codim_X(Z)\>2$,  and $g:X'\to X$ a log resolution of $X_{\sing}\cup Z$ such that $X'$ is a K\"ahler manifold. Then by \cite[Theorem 3.17(ii)]{Bou04}, there is a K\"ahler current $T$ with analytic singularities in $g^*\alpha$ such that $E_+(T)=E_{nK}(g^*\alpha)$. Let $h:X''\to X'$ be a resolution of singularities of $T$ such that $h^*T=\phi+D$, where $D\>0$ is an effective $\mbR$-divisor and $\phi$ is a smooth form (see \cite[Definition 2.5.1]{Bou04}). Since by Lemma \ref{lem:non-kahler-locus}, $E_{nK}(g^*\alpha)=g^{-1}(E_{nK}(\alpha))\cup g^{-1}(g(\Ex(g))$ and $E_{nK}(\alpha)\subset Z$, it follows that $D$ is $g\circ h$-exceptional. Let $\omega'$ be a K\"ahler form on $X'$ such that $T\geq \ve\omega'$ for some $\ve>0$. Then $\phi\geq \ve\mu^*\omega'$, and thus by \cite[Lemma 2.9]{Bou02ar} we see that there is an effective $h$-exceptional $\mbR$-divisor $E\>0$ such that $\phi -\delta E$ is cohomologous to a K\"ahler form for $0<\delta \ll 1$. Then $(g \circ h) ^*\alpha \equiv \omega '' +G$, where $\omega ''$ is a K\"ahler form and $G\geq 0$ is $g \circ h$-exceptional effective $\mbR$-divisor.  Finally, set $f:=g\circ h$, and $\omega_Y:=\omega''$. Then we have $f_*\omega_Y=f_*(\omega_Y+G)=(g\circ h)_*(\omega''+G)=\alpha$. This completes our proof.  
 \end{proof}




\subsection{Adjunction for Generalized Pairs}\label{subsec:adjunction} In this subsection we recall the adjunction for usual pairs and define adjunction for generalized pairs.\\
\begin{itemize}
    \item Recall that, if $(X, S+B)$ is a log canonical pair such that $\lrd S+B\rrd=S$ is irreducible and $S^\nu\to S$ is the normalization, then by adjunction (see \cite[Chapter 4]{Kol13} there is an effective $\mbQ$-divisor ${\rm Diff}_{S^\nu}(B)\>0$ on $S^\nu$, called the \emph{different}, such that 
\[  
(K_X+S+B)|_{S^\nu}=K_{S^\nu}+{\rm Diff}_{S^\nu}(B).
\]
\end{itemize}
Now let $(X, B+\bbeta)$ be a glc pair such that $X$ is in Fujiki's class $\mcC$ and has rational singularities. Assume that $K_X$, $B$ and every component of $\lrd B\rrd$ is $\mbQ$-Cartier ($K_X$ $\mbQ$-Cartier means that there is a $m\in\mbZ^+$ such that the reflexive sheaf $(\omega_X^{\otimes m})^{**}$ is a line bundle). 
Fix a component $S$ of $\lrd B\rrd$ and let $f:X'\to X$ be a log resolution of the glc pair $(X, B+\bbeta)$, $S'$ is the strict transform of $S$ and $K_{X'}+B'+\bbeta_{X'}=f^*(K_X+B+\bbeta_X)$. Since $K_X+B$ is $\mbQ$-Cartier, it follows that $\bbeta_X$ is a closed positive $(1,1)$ current with local potentials on $X$. Let $S^\nu\to S$ be the normalization morphism. We would like to define \[K_{S^\nu}+B_{S^\nu}+\bbeta_{S^\nu}=(K_X+B+\bbeta_X)|_{S^\nu}={f|_{S'}}_*(K_{S'}+(B'-S')|_{S'}+\bbeta_{X'}|_{S'}).\] 
However, the difficulty here is that $S'$ could be contained in the pluripolar locus of the local potentials of $\bbeta_{X'}$, in which case $\bbeta_{X'}|_{S'}$ is not well defined. To remedy this, we will replace $\bbeta_{X'}$ by a representative in its class $[\bbeta_{X'}]\in H^{1,1}_{\BC}(X')$ that does not contain $S'$ in its pluripolar locus.


Let $\alpha_{X'}$ be a smooth $(1,1)$ form on $X'$ such that $[\alpha_{X'}]=[\bbeta_{X'}]$. Then by \cite[Lemma 8.7]{DHP22} it follows that $\alpha_X=f_*\alpha_{X'}$ is a $(1,1)$ current on $X$ such that $[\alpha_X]=[\bbeta_X]$ in $H^{1,1}_{\BC}(X)$, and $\alpha_{X'}$ defines a b-$(1,1)$ current $\boldsymbol{\alpha}:=\overline{\alpha_{X'}}$ on $X$. From the negativity lemma we have $f^*\bbeta_X=\bbeta_{X'}+E$ and $f^*\alpha_X=\alpha_{X'}+F$, where $E$ and $F$ are effective $f$-exceptional $\mbR$-divisors. Then   $E-F\num 0$, and thus applying the negativity lemma again we get that $E=F$. In particular, for any prime Weil divisor $P$ over $X$ we have $a(P, X, B+\boldsymbol{\alpha}_X)=a(P, X, B+\bbeta_X)$.  Then we define 
\begin{eqnarray*}
    K_{S^\nu}+B_{S^\nu}+\bbeta_{S^\nu}:=(K_X+B+\bbeta_X)|_{S^\nu} &\num & (K_X+B)|_{S^\nu}+\boldsymbol{\alpha}_X|_{S^\nu}\\
                                                                &= &{f|_{S'}}_*(K_{S'}+(B'-S')|_{S'}+\boldsymbol{\alpha}_{X'}|_{S'}).
\end{eqnarray*}
From Proposition 4.5 and 4.7 of \cite{Kol13} it follows that $B_{S^\nu}$ is a boundary divisor.\\

\begin{remark}
 Note that $K_{X'}+B'-E=f^*(K_X+B)$, where $0\<E=f^*\bbeta _X-\bbeta _{X'}$ is an effective $f$-exceptional $\mbR$-divisor. It then follows easily from the definition above that if $\phi=f|_{S'} :S'\to S^\nu$ is the induced morphism, then 
 \[B_{S^\nu}={\rm Diff}_{S^\nu}(B)+\phi _*(E|_{S'})\geq {\rm Diff}_{S^\nu}(B).\]
\end{remark}



\subsection{Varieties in Fujiki's class $\mcC$} Recall that a compact analytic variety $X$ is said to be in Fujiki's class $\mcC$ if $X$ is bimeromorphic to a compact K\"ahler manifold. In what follows, we will study some basic properties of these varieties which will be useful in running the MMP.

\begin{lemma}\label{lem:pullback-nef}
Let $X$ be a normal compact variety of dimension $3$ in Fujiki's class $\mcC$ and $f:Y\to X$ is a proper bimeromorphic morphism from a normal variety $Y$. Then a class $\alpha\in H^{1,1}_{\BC}(X)$ is nef if and only if $f^*\alpha$ in nef on $Y$
\end{lemma}
\begin{proof}
    This is \cite[Lemma 3.13]{HP16}.
\end{proof}

\begin{lemma}\label{lem:f-adjoint-nef}
 Let $(X, B)$ be a $\mbQ$-factorial dlt pair, where $X$ is a compact variety of dimension $\<3$ in Fujiki's class $\mcC$. If $K_X+B$ is pseudo-effective, then $K_X+B$ is nef if and only if $(K_X+B)\cdot C\>0$ for every compact curve $C\subset X$.
\end{lemma}

\begin{proof}
We will show the proof for the dimension $3$ case and leave the surface case as an easy exercise to the reader. The only if part follows easily by passing to a resolution of singularities $f:X'\to X$ such that $X'$ is a K\"ahler manifold. Conversely, assume that $(K_X+B)\cdot C\>0$ for all curves $C\subset X$. Then by \cite[Remark 2.36]{DHP22}, there is a surface $S\subset X$ such that $(K_X+B)|_S$ is pseudo-effective. Let $K_X+B\num \sum\lambda_iS_i+\beta$ be the Boucksom-Zariski decomposition of $K_X+B$ as in Definition \ref{def:mod-nef}. Then by a similar argument as in the proof of \cite[Lemma 4.1]{HP16} it follows that $K_{S^\nu}$ is not pseudo-effective, where where $\nu:S^\nu\to S$ is the normalization. Let $\tilde S\to S^\nu$ be the minimal resolution of $S^\nu$. Then $K_{\tilde S}$ is also not pseudo-effective; then from classification (see \cite[Table 10, page 244]{BHPV04}) it follows that $\kappa(\tilde S)=-\infty$. Since $\tilde S$ is in Fujiki's class $\mcC$, by \cite[Theorem 1.3]{Fuj21} it follows that $\tilde S$ is projective, and thus $S$ is Moishezon. Then we arrive at a contradiction by a similar argument as in \cite[Corollary 4.2]{HP16}. 
\end{proof}


\begin{lemma}\label{lem:fujiki-to-kahler-surface}
    Let $(S, B+\bbeta )$ be a gdlt pair, where $S$ is a compact analytic surface in Fujiki's class $\mcC$. Then $S$ is a K\"ahler surface with $\mbQ$-factorial rational singularities. 
\end{lemma}

\begin{proof}
Since $(S, B+\bbeta)$ is gdlt, by \cite[Lemma 2.20]{DHY23}, $S$ has $\mbQ$-factorial rational singularities.

Now let $\nu :S'\to S$ be a resolution of singularities of $S$ such that $S'$ is a K\"ahler manifold of dimension $2$. Let $\omega '$ be a K\"ahler class on $S'$. By the Hodge index theorem, the intersection matrix of the exceptional curves of $\nu$ is negative definite. Thus there is a unique $\nu$-exceptional $\mbR$-divisor $E$ such that $\omega'+E\num_S 0$. Since $S$ has rational singularities and it belongs to Fujiki's class $\mcC$, by \cite[Lemma 3.3]{HP16} there is a class $\omega\in H^{1,1}_{\BC}(S)$ such that $\omega'+E=\nu^*\omega$. From the negativity lemma it follows that $E\>0$. Thus $\omega=\nu_*(\omega'+E)=\nu_*\omega'$ is a big class on $S$. 

If $C\subset S $ is any (compact) curve and $C':=\nu ^{-1}_*C$, then $\omega \cdot C=(\omega '+E)\cdot C'>0$. Then by \cite[Theorem 2.36 and Remark 2.37]{DHP22}, $\omega$ is nef, and by \cite[Theorem 2.30]{DHP22}, $\omega $ is a  K\"ahler class.
\end{proof}


\begin{theorem}[Termination of Flips]\label{thm:termination}
    Let $(X, B)$ be a $\mbQ$-factorial dlt pair, where $X$ is a compact analytic $3$-fold in Fujiki's class $\mcC$. Then every sequence of $K_X+B$-flip terminates.
\end{theorem}

\begin{proof}
First we note that Fujino's proof of Special termination (see \cite{Fuj07}) holds here as the log MMP for surfaces in Fujiki's class $\mcC$ is known again due to Fujino, see \cite{Fuj21}. So we may assume hat $(X, B)$ is a klt pair. Next we claim that the proof presented in \cite[Theorem 3.3]{DO23} (which is based on the proof of \cite{Kaw92}) holds in our case without any change. This is because even though $X$ is assumed to be K\"ahler in \cite{DO23}, that property was never used in the proof, and all the necessary MMPs used in that proof are relative MMPs for projective morphisms which are known to exist due to \cite{Nak87} and more recently by \cite{Fuj22} and \cite{DHP22}. 
\end{proof}












\section{Blowing down analytic spaces to points}\label{sec:point-contractions}
Recall the following result, generalizing \cite{Gra62}, that is claimed in \cite[Proposition 7.4]{HP16}.
\begin{claim}\label{c-HP16} Let $X$ be a normal compact complex space, and let $S$
be a prime divisor that is $\Q$-Cartier of Cartier index $m$. Suppose that $S$
admits a morphism with connected fibres $f: S \to T$ such that $\mathcal O _S (-mS)$
is $f$-ample. Then there exists a bimeromorphic morphism $\varphi :X\to Y$  to
a normal compact complex space $Y$ such that $\varphi|_S = f$ and $\varphi|_{X\setminus S}$ is an
isomorphism onto $Y \setminus T$.
\end{claim}
The above claim contradicts \cite[Proposition 3]{Fuj74}, and the proof seems to confuse the divisor $S$ with its thickenings (a similar issue occurs in \cite{CHP16}). To be more precise the strategy of \cite[Proposition 7.4]{HP16} is based on \cite{AT84} which also appears to contain an error. One should instead use the corresponding result in \cite[Theorem 2]{Fuj74} which states:
\begin{theorem}\cite[Theorem 2]{Fuj74}\label{t-fuj74}
Let $X$ be a reduced complex space, $A$ an effective Cartier
divisor on $X$ and $f: A\to A'$ a proper surjective morphism of $A$ onto another
complex space $A'$.  Assume that the following conditions hold:
\begin{enumerate}
    \item $\mathcal O _A(-A)$ is $f$-ample, and
    \item $R^1f_*\OO _A(-kA)=0$ for all $k>0$. 
\end{enumerate}
Then there exists a blowing down $F:X\to X'$ such that $F_*\mathcal S_{X,A,f}=\OO _{X'}$ where $\mathcal S_{X,A,f}$ is the coherent sheaf defined by the following short exact sequence
\[ 0\to \mathcal S_{X,A,f}\to \OO _X \to \OO _A/f^{-1}\OO _{A'}\to 0.\] 
\end{theorem}
The proof of \cite[Proposition 7.4]{HP16} begins by checking that $R^1f_* \OO _S(-km'S)=0$ for $k>0$, and given $\mathcal I =\OO _X(-m'S)$ it is claimed that 
\[R^1f_* (\mathcal I^k/\mathcal I^{k+1})\cong R^1f_*\OO _S(-km'S)\] however the right hand side should presumably be replaced by $R^1f_*\OO _{kS}(-km'S)$ (note that the prime divisor $S$ is not assumed to be Cartier in \cite{HP16}, while \cite[Theorem 2]{Fuj74} requires Cartier divisors).
When $\dim A'=0$, one can however deduce the following result of Grauert \cite{Gra62}.
\begin{lemma}\label{l-cont}
Let $X$ be a normal analytic variety and $S\subset X$ a prime Weil divisor such that $S$ is compact.  Assume that $S$ is $\Q$-Cartier of Cartier index $m>0$ such that $\mcO_S(-mS|_S)$ is ample. Then there exists a proper bimeromorphic morphism $F:X\to Y$ to a normal analytic variety $Y$ such that $F(S)=\{p_0\}\subset Y$ is a point and $F|_{X\setminus S}:X\setminus S\to Y\setminus \{p_0\}$ is an isomorphism. 
\end{lemma}
\begin{proof}
Since $\mcO_S(-mS|_S)$ is ample, from \cite[Proposition 1.2.16(i)]{Laz04a} it follows that  $\mcO_{mS}(-mS)$ is ample. Thus by Serre's vanishing theorem, there is a positive integer $k_0>0$ such that $H^i(mS, \OO _{mS}(-kmS))=0$ for all $i>0$ and $k\geq k_0$.
     Thus from the short exact sequences (for $j\in \mbZ^{\>0}$)
    \[0\to \OO _{mS}(-(k+j)mS)\to \OO _{(j+1)mS}(-kmS)\to \OO _{jmS}(-kmS) \to 0\]
    it follows that $H^i((j+1)mS, \OO _{(j+1)mS}(-kmS))=0$ for all $i>0$, $j> 0$ and $k\geq k_0$.
    In particular, $H^i(k_0mS, \OO _{k_0mS}(-jk_0mS))=0$ for all $i,j>0$.
    Replacing $m$ by $k_0m$, we may assume that $H^i(mS, \OO _{mS}(-jmS))=0$ for all $i,j>0$.
   
    
    We let $A=mS$. Then $A$ is Cartier and $\mcO_A(-A|_A)$ is ample since $\mcO_S(-mS|_S)$ ample (see \cite[Proposition 1.2.16(i)]{Laz04a}). Note that $A$ is a projective scheme (of finite type) over $\mbC$. Let $f:A\to \Spec \mbC$ be the structure morphism. Then we have $R^1f_*\OO _A(-kA)=H^1(mS, \OO _{mS}(-kmS))=0$ for all $k>0$, and thus by Theorem \ref{t-fuj74} there is a blowing down $F:X\to X'$ such that $F(A)=\pt$ and $F$ satisfies the other required properties. Since $\Supp(S)=\Supp (A)$, we are done.
    
   
   
\end{proof}~\\

As a corollary we will show below that if $\alpha$ is a nef $(1,1)$ class and $S:=\Null(\alpha)$ is an irreducible surface such that $\alpha|_{S^\nu}\num 0$, where $S^\nu\to S$ is the normalization, then we can contract $S$. 





%As a corollary to Lemma \ref{l-cont} we prove the existence of 3-fold divisorial contractions to a point.
\begin{corollary}\label{c-1}
Let $X$ be a normal $\Q$-factorial compact K\"ahler variety of arbitrary dimension and 
$\alpha\in H^{1,1}_{\rm BC}(X)$ a nef and big class such that the following hold:
\begin{enumerate}
\item $\alpha ^\perp$ defines an extremal ray $R$ in $\overline{\rm NA}(X)$, 
\item the curves in $R$ cover a prime Weil divisor $S\subset X$, and 
\item $\alpha |_{S^\nu}\equiv 0$, where $S^\nu\to S$ is the normalization morphism. 
\end{enumerate}
Then there exists a proper bimeromorphic contraction  $f:X\to Y$ which contracts $S$ to a point and is the identity morphism on the complement of $S$.


\end{corollary}



\begin{proof}
Since $\alpha$ is big, we may write $\alpha=\gamma +\omega$, where $\omega $ is a K\"ahler class and $\gamma$ is a big class. Let $\gamma=Z+\beta$ be the Boucksom-Zariski decomposition of $\gamma$ as defined above, where $Z$ is an $\mbR$-divisor and $\beta$ is a modified nef class. Then $\beta |_{S^\nu}$ is pseudo-effective.  If $S$ is not contained in the support of $Z$, then $\alpha |_{S^\nu} \equiv (Z+\beta+\omega)|_{S^\nu}$ is big, which is impossible as $ \alpha |_{S^\nu} \equiv 0$.
Thus $S$ is  contained in the support of $Z$. We have $0= \alpha \cdot R = (Z+\beta+\omega)\cdot R$ and so $T\cdot R<0$, where $T=\beta$ or $T$ is a component  of $Z$. Note that $\NA(X)=\NA(X)_{T\>0}+\NA(X)_{T<0}$, so from a standard argument (e.g. see the proof of \cite[Claim 3.25]{DHY23}) it follows that $\alpha -\delta T$ is positive on $\overline{\rm NA}(X)$ for some $\delta>0$, and hence K\"ahler. But then $(\alpha -\delta T)|_{S^\nu}\equiv -\delta T|_{S^\nu}$ K\"ahler. 
Since $\beta$ is modified nef, $\beta|_{S^\nu}$ is pseudo-effective and so $T$ is a component  of $Z$.  
Thus $T=S$, as otherwise $-T|_{S^\nu}$ is an anti-effective class on $S^\nu$ and thus cannot be K\"ahler. In particular,  $-S|_S$ is K\"ahler, and hence $\mcO_S(-mS|S)$ is an ample line bundle, where $m$ is the Cartier index of $S$ on $X$. Then by Lemma \ref{l-cont} there is a bimeromorphic contraction $f:X\to Y$ such that $f(S)=\{p_0\}$ is a point and $f|_{X\setminus S}: X\setminus S\to Y\setminus \{p_0\}$ is an isomorphism.

\end{proof}


 














A second consequence of Lemma \ref{l-cont} is the existence of flips for terminal 3-folds (or more generally lc 3-folds with isolated rational singularities).

\begin{proposition}\label{p-1}
Let $X$ be a normal compact K\"ahler variety of arbitrary dimension with isolated rational singularities and $\alpha$ a nef and big $(1,1)$ class on $X$. If  $\Null(\alpha)$ is a pure $1$-dimensional analytic subspace of $X$, then there is a proper bimeromorphic morphism $f:X\to Y$ such that $\Ex(f)=\Null(\alpha)$ and $f(\Null(\alpha))$ is a finite set of points.

\end{proposition}
\begin{proof}
   We will follow the arguments in the proof of \cite[Theorem 7.12]{HP16}. Let $\nu :X'\to X$ be a log resolution of $X$ and $\Null(\alpha)$, which is an isomorphism on the complement of $X_\textsubscript{sing}\cup {\rm Null}(\alpha)$. Then $\nu ^*\alpha $ is also nef and big and hence, by \cite[Theorem 1.1]{CT15}  and Lemma \ref{lem:null-locus} we have
   \[ E_{nK}(\nu ^*\alpha )=\Null(\nu ^*\alpha)=\Ex(\nu)\cup \nu^{-1}(\Null(\alpha)).\]
Since $E_{nK}(\alpha )\subset \nu (E_{nK}(\nu ^*\alpha ))$ by Lemma \ref{lem:non-kahler-locus}, $E_{nK}(\alpha)$ doesn't not contain any divisor of $X$. By \cite[Theorem 3.17(ii)]{Bou04}, there exists a K\"ahler current $T$ in $\nu ^*\alpha$ with analytic singularities such that
$E_+(T ) = E_{nK} (\nu ^* \alpha )$. Let $\mu :X''\to X'$ be a resolution of the singularities of $T$ so that $\mu ^*T=\phi +D$, where $D\>0$ is an effective $\R$-divisor and $\phi$ is a smooth form (see \cite[Definition 2.5.1]{Bou04}); note that $D$ is $f\circ\nu$-exceptional.
% \footnote{Om: This doesn't seem to follow without the equality $E_{nK}=\nu(E_{nK}(\nu^*\alpha)$, and without $D$ being exceptional rest of our argument will not work!}
Let $\omega'$ be a K\"ahler form on $X'$ such that $T\geq \ve\omega'$ for some $\ve>0$. Then $\phi\geq \ve\mu^*\omega'$, and thus by \cite[Lemma 2.9]{Bou02ar} we see that there is an effective $\mu$-exceptional $\mbR$-divisor $E\>0$ such that $\phi -\delta E$ is cohomologous to a K\"ahler form for $0<\delta \ll 1$. Then $(\nu \circ \mu) ^*\alpha \equiv \omega '' +G$, where $\omega ''$ is a K\"ahler form and $G\geq 0$ is $f \circ \mu$-exceptional effective $\mbR$-divisor.
Let $C^{\nu}\to C$ be the normalization of $C$, a connected component of $\mathcal N:={\rm Null}(\alpha)$. Then $\alpha |_{C^\nu}\equiv 0$ (since $C$ is a curve) and so if $D:=(\nu \circ \mu) ^{-1}(C)$, then $D$ is a reduced effective divisor contained in the support of $G$ and
$-G|_D\equiv \omega ''|_D$ is ample. After perturbing $G$ (without changing its support), we obtain an effective $\mbQ$-divisor $G'$ such that $-G'|_D$ is ample. From our construction above it follows that $(\nu\circ\mu)(\Ex(\nu\circ\mu))=X_{\sing}\cup\Null(\alpha)$. Thus the components of $G$ which do not map into $\Null(\alpha)$ are contracted to (the singular) points of $X$, and hence $-G|_{{\rm Supp}(G)}$ is ample. So there is a contraction $X''\to Y$ which contracts ${\rm Supp}(G)$ to a finite set of points and is an isomorphism on the complement of ${\rm Supp}(G)$. By the rigidity lemma (see eg. \cite[Lemma 4.1.13]{BS95}), we obtain the required morphism $X\to Y$.

\end{proof}


\begin{corollary}
    Let $(X,B)$ be a compact K\"ahler $3$-fold terminal pair. Let $\omega$ be a K\"ahler class on $X$ and $\alpha=K_X+B+\omega$ a nef but not K\"aher class on $X$ such that $\Null(\alpha)$ is a pure $1$-dimensional analytic subspace of $X$ and $\overline{\rm NA}(X)\cap \alpha ^\perp =R$ is an extremal ray. Then the flipping contraction $f:X\to Z$ and the corresponding flip $f^+:X^+\to Z$ both exist.
\end{corollary}
\begin{proof}
    As $(X,B)$ is  terminal, $X$ has only isolated singularities.
    Thus Proposition \ref{p-1} applies and hence the flipping contraction $f:X\to Z$ exists.
    By \cite[Theorem 1.8]{Fuj22} or \cite[Theorem 1.3]{DHP22}, we can construct the log canonical model of $(X,B)$ locally over $Z$ and then glue these together to obtain the flip $f^+:X^+\to Z$.
\end{proof}
In order to complete the 3-fold MMP for klt pairs, it is necessary to construct divisorial contractions where $n(\alpha )=1$, i.e. divisorial contractions to a curve,  and flipping contractions in the non-terminal case.
Divisorial contractions to a curve were constructed in \cite{DH20} and flipping contractions are constructed in \cite{CHP16}. Note however that the proof of \cite{CHP16} is somewhat technical as it involves the use of ample sheaves (which are not necessarily vector bundles). In what follows we will give a unified approach to the construction of flipping contractions and divisorial contractions to a curve. This approach relies on first proving that the pl-extremal contractions exist. Our approach seems to be well suited for higher dimensional arguments and settles several important higher dimensional cases. 

%%%%%%%%%%%%%%%%%%%%%%%%%%%%%%%%%%%%%%%%%
\section{The Main Contraction Theorems}\label{sec:main-contractions}
The goal of this section is to prove the contraction Theorems \ref{t-pl} and \ref{thm:contractions}.
\subsection{ Existence of pl-contractions and flips for generalized pairs} In this subsection we will prove Theorem \ref{t-pl}. First we need some preparatory results. The following theorem is a version of relative Kawamata-Viehweg vanishing theorem which holds for generalized pairs. 
\begin{lemma}\label{l-wdvan} Let $(X,B+\bbeta)$ be a  gklt pair, $g:X\to T$ a projective morphism where $T$ is relatively compact and $D$ a $\Q$-Cartier Weil divisor such that $D-(K_X+B+\bbeta_X)$ is nef and big over $T$. Then $R^ig_* \OO _X(D)=0$ for $i>0$.
\end{lemma}
\begin{proof}
Let $\nu :(X',B'+\bbeta_{X'})\to (X,B+\bbeta_X) $ be a log resolution of the gklt pair $(X,B+\bbeta)$ such that $K_{X'}+B'+\bbeta_{X'}=\nu^*(K_X+B+\bbeta_X)$. 
We have $\lceil \nu ^*D\rceil=\nu ^*D+\{-\nu^*D\}$, and we let  $E:=\lfloor B'+\{-\nu ^*D\}\rfloor $, where $\{\cdot\}$ denotes the fractional part. 

We make the following claim.
\begin{claim}\label{clm:divisor-comparison}
  $\lceil \nu ^*D\rceil -E\geq \lfloor \nu ^*D\rfloor $.  
\end{claim}
  
\begin{proof}[Proof of Claim \ref{clm:divisor-comparison}]
Let $P$ be any prime Weil divisor on $X'$. First assume that $P$ is not contained in the support of $\{-\nu ^*D\}$. Then ${\mult}_P(-E)={\mult}_P(-\lfloor B'\rfloor )\geq 0$ and ${\mult}_P(\lceil \nu ^*D\rceil) =\mult_P(\nu^*D)={\mult}_P(\lfloor \nu ^*D\rfloor )$, and the claim follows.  Suppose now that $P$ is contained in the support of $\{-\nu ^*D\}$. Then $\mult_P(\{\nu^*D\})=(1-\mult_P(\{-\nu^*D\}))>0$, so that ${\rm mult}_P(\lceil \nu ^*D\rceil) = {\rm mult}_P(\lfloor \nu ^*D\rfloor )+1$. Since  ${\rm mult}_P(-E)={\rm mult}_P(-\lfloor B'+\{-\nu ^*D\}\rfloor )\geq -1$, the required inequality is clear.  
\end{proof}  


Now recall that $(X, B+\bbeta)$ is a gklt pair and $D$ is a Weil divisor. Therefore $E=\lrd B'+\{-\nu^*D\}\rrd$  is $\nu$-exceptional, $\nu _* (\lceil \nu ^*D\rceil-E)=D$ and $\nu_*\lrd \nu^*D\rrd =D$. Since $\nu$ is bimeromorphic, the natural morphism $\nu_*\mcO_{X'}(\lru\nu^*D\rru-E)\to \mcO_X(\nu_*(\lru\nu^*D\rru-E))=\mcO_X(D)$ is injective, and $\nu_*\mcO_{X'}(\lrd\nu^*D\rrd)=\mcO_X(D)$ (see \cite[Lemma 3.2]{Das21}).  By Claim \ref{clm:divisor-comparison} it follows that
\[\OO _{X}(D)\supset \nu _* \OO _{X'}(\lceil \nu ^*D\rceil -E)\supset \nu _* \OO _{X'}(\lfloor \nu ^*D\rfloor)=\OO _X(D).\] 
Therefore we have
\begin{equation}\label{eqn:kv-isomorphism}
    \nu_*\mcO_{X'}(\lru\nu^*D\rru-E)=\mcO_X(D).
\end{equation}
 Now observe that 
 \[\lceil \nu ^*D\rceil -E= K_{X'}+B^*+\bbeta _{X'}+\nu ^*(D-(K_X+B+\bbeta_X)),\]
where $B^*:=B'+\{-\nu ^*D\}-E=\{B'+\{-\nu ^*D\}\}$. But then $(X',B^*)$ is klt and the $\R$-divisor \[\lceil \nu ^*D\rceil -E-(K_{X'}+B^*)\equiv \bbeta _{X'}+\nu ^*(D-(K_X+B+\bbeta _X))\] is  nef and big over $T$, and so, by the relative Kawamata-Viehweg vanishing theorem \cite[Theorem 3.7]{Nak87} (see also \cite[Theorem 5.2]{Fuj22}) for all $i>0$ we have 
\[R^i\nu _* \OO _{X'}(\lceil \nu ^*D\rceil -E)=0,\ \qquad R^i(g\circ \nu )_* \OO _{X'}(\lceil \nu ^*D\rceil -E)=0.\]
By a standard spectral sequence argument and \eqref{eqn:kv-isomorphism} it follows that 
\[R^ig_*\OO _X(D)=R^ig_* (\nu _* \OO _{X'}(\lceil \nu ^*D\rceil -E))=0\qquad \mbox{ for all } \ i>0.\]


\end{proof}








\begin{lemma}\label{lem:normality-of-plt-center}
    Let $(X, S+B+\bbeta)$ be a gplt pair such that $\lrd S+B\rrd=S$ is irreducible, where $X$ is a compact analytic variety. Then $S$ is normal and $(S, B_S+\bbeta_S)$ is gklt, where $K_S+B+\bbeta_S=(K_X+S+B+\bbeta_X)|_S$ is defined by adjunction as in Subsection \ref{subsec:adjunction}.
\end{lemma}

\begin{proof}

    Let $f:X'\to X$ be a log resolution of the generalized pair $(X, B+\bbeta)$ and 
\[ 
K_{X'}+S'+B'+\bbeta_{X'}=f^*(K_X+S+B+\bbeta_X)+E'
\]
where $S'=f^{-1}_*S, B'\>0, E'\>0, f_*B'=B, f_*E'=0$ and $S'+B'$ and $E'$ do not share any common component.

Then $-S'+\lru E'\rru \num_f K_{X'}+B'+\{-E'\}+\bbeta_{X'}$ and $(X', B'+\{-E'\})$ is a klt pair. Since $[\bbeta_{X'}]\in H^{1,1}_{\BC}(X')$ is a nef class,

from Lemma \ref{l-wdvan} it follows that $R^if_*\mcO_{X'}(-S'+\lru E'\rru)=0$ for all $i>0$. Now consider the following exact sequence 
\begin{equation*}
\xymatrixcolsep{3pc}\xymatrix{
0\ar[r] & \mcO_{X'}(\lru E'\rru -S')\ar[r] & \mcO_{X'}(\lru E'\rru)\ar[r] & \mcO_{S'}(\lru E'\rru|_{S'})\ar[r] & 0.
}
\end{equation*}
Since $R^1f_*\mcO_{X'}(-S'+\lru E'\rru)=0$, from the long exact sequence of cohomology we get the surjection 
\[ 
f_*\mcO_{X'}(\lru E'\rru)\surjective f_*\mcO_{S'}(\lru E'\rru|_{S'}).
\]
Since $\lru E'\rru|_{S'}$ is an effective divisor, this map factors as
\[
\mcO_{X}=f_*\mcO_{X'}(\lru E'\rru)\to \mcO_{S}\injective \nu_*\mcO_{S^\nu}\injective f_*\mcO_{S'}(\lru E'\rru|_{S'}),
\]
where $\nu:S^\nu\to S$ is the normalization morphism. Thus we have $\mcO_{S}\cong\nu_*\mcO_{S^\nu}$, and hence $S$ is normal.\\


For the second part by contradiction assume that there is a prime Weil divisor $F$ over $S$ such that $a(F, S, B_S+\bbeta_S)\<-1$. Passing to a sufficiently high resolution $f:X'\to X$ we may assume that we are in the setup of Subsection \ref{subsec:adjunction} and that there is a $f$-exceptional divisor $E$ on $X'$ such that $E\cap S'=F$ and $f(E)=f|_{S'}(F)$. It then follows that $a(E, X, S+B+\bbeta_X)=a(F, S, B_S+\bbeta_S)\<-1$, a contradiction, since $(X, S+B+\bbeta)$ is gplt.   

\end{proof}~\\



\begin{lemma}\label{lem:CM-sheaves}
    Let $(X, B)$ be a klt pair, where $X$ is an analytic variety and $D$ a $\mbQ$-Cartier Weil divisor. Then $\mcO_X(D)$ is Cohen-Macaulay.
\end{lemma}

\begin{proof}
    The proof follows form \cite[Corollary 5.25]{KM98}. We explain the necessary changes needed for analytic variety. Since the question is local on the base, replacing $X$ by a relative compact Stein open  subset we may assume that $X$ is relatively compact Stein space. Let $m>0$ be the Cartier index of $D$, then $\mcO_X(-mD)$ is a line bundle on $X$. Since $X$ is Stein, $\mcO_X(-mD)$ is generated by its global sections. Using $\mcO_X(-mD)$ in place of $\mcO_X(mL-mD)$ in the proof of \cite[Corollary 5.25]{KM98}, and using \cite[Theorem 2.21]{DHP22} for the necessary Bertini-type theorems, the arguments of \cite[Corollary 5.25]{KM98} work here.  
    
    
\end{proof}





Proposition \ref{p-hw} below is an analytic version of a key result from \cite{HW19} (see also \cite{BK23}), and an important step the proof of Theorem \ref{t-pl}. First we  fix some notation.

\begin{definition}\label{def:Q-line-bundle}
Let $X$ be a normal analytic variety.
    \begin{enumerate}
   
    \item A non-empty open subset $U\subset X$ called \emph{big}, if $X\setminus U$ is a countable union of locally closed (analytic) subsets of $X$ of codimension at least $2$.

    \item A sheaf $\mcL$ on $X$  is called a $\Q$-{\emph {line bundle}} $\mathcal L$, if $\mathcal L$ is a reflexive sheaf of rank $1$ and $(\mathcal L^{\otimes m})^{**}$ is locally free for some $m>0$. 
Note that if $D$ is a $\Q$-Cartier Weil divisor, then $\mathcal L=\OO _X(D)$ is a $\Q$-line bundle. On the other hand, if $X$ is Stein, then for any $\Q$-line bundle $\mathcal L$ there exists a $\Q$-Cartier Weil divisor $D$ such that $\mathcal L=\OO _X(D)$. Moreover, 
If $\mathcal L$ is a $\Q$-line bundle and $G$ is a Weil divisor, then we define  $\mathcal L(D):=(\mathcal L \otimes \OO _X(G))^{**}$ and $\mathcal L(K_X+D):=(\mathcal L \otimes \omega _X\otimes \OO _X(G))^{**}$.
    \end{enumerate}  
\end{definition}


\begin{proposition}\label{p-hw} 
    Suppose that $(X,S)$ is a relatively compact plt pair, $S$ is a $\Q$-Cartier Weil divisor and $\Delta _{\rm diff}:={\rm Diff}_S(0)$ is the different.
    Then for any $\Q$-line bundle $\mathcal L$ %$\Q$-Cartier Weil divisor $D$
    on $X$, there exists an effective $\Q$-divisor $0\<\Delta _S \leq \Delta _{\rm diff}$, a positive integer $m>0$,  
and a  reflexive rank 1 sheaf $\mathcal E$ on $S$  such that the following sequence
\begin{equation}\label{eqn:reflexive-exact}
0\to \mathcal L(K_X)\to \mathcal L(K_X+S)\to \mathcal E \to 0
\end{equation}
is exact, and there is a smooth big open subset $S^0\subset S$ such that $\mathcal E|_{S^0}$ is locally free, $m\Delta _S$ is a $\mbZ$-divisor, %$\mathcal O _X(mD)|_{S^0}$ 
$(\mathcal L^{\otimes m})|_{S^0}$ is locally free,  and \[(\mathcal E|_{S^0}) ^{\otimes m}\cong (\mathcal L^{\otimes m})|_{S^0}\otimes \mathcal O _{S_0}(m(K_S+\Delta _S)).\] 
\end{proposition}

\begin{proof} This proof follows along the lines of \cite[Proposition 3.1]{HW19}. We include the details so that we can highlight the changes necessary for the analytic context.
    Note that since $(X,S) $ is plt,  $S$ is normal (see Lemma \ref{lem:normality-of-plt-center}).
      
    Let $\mu:X'\to X$ be a log resolution of $(X,S)$ which is given by a sequence of blow ups along smooth centers, $S'=\mu ^{-1}_*S$, and $\Delta _{X'}$ defined by the equation
    \[K_{X'}+S'+\Delta _{X'}=\mu ^*(K_X+S). \]
    

    Let $X=\cup U_i$ be a cover by Stein open subsets. We may write $\mathcal L |_{U_i}=\OO _{U_i}(D_i)$ where $D_i$ is Weil and $\Q$-Cartier. We define $U'_i:=\mu ^{-1}(U_i)$, $\mu _i:=\mu |_{U'_i}$, $S_i:=S\cap U_i$, $S'_i:=S'\cap U'_i$, $\nu _i:=\mu _i|_{S'_i}$ and $\Delta_{S_i}:=\Delta _{\rm diff}|_{S_i}-\nu _{i,*} (\{-\mu_i^*D_i\}|_{S'_i})$.
    Note that the support of $\{-\mu_i^*D_i\}$ does not contain $S_i$ and if $\bar D_i\sim D_i$, then $\{-\mu_i^*D_i\}=\{-\mu_i^*\bar D_i\}$. Therefore, the definition of $\Delta_{S_i}$ is independent of our choice of $D_i$. It then follows that, on $U_{i,j}:=U_i\cap U_j$, we have $D_i|_{U_{i,j}}\sim D_j|_{U_{i,j}}$ and hence $\Delta_{S_i}|_{S_{i,j}}=\Delta_{S_j}|_{S_{i,j}}$, where $S_{i,j}:=S_i\cap S_j$.
    Therefore, the $\Delta_{S_i}$'s glue together to define a unique $\mbQ$-divisor $\Delta_S$ on $S$ such that $\Delta_S |_{S_i}=\Delta_{S_i}$.
    
    Since all remaining claims are local on $X$, we may assume that $X$ is Stein. In particular, we may assume that $\mathcal L=\OO _X(D)$ for some $\Q$-Cartier Weil divisor $D$. We then have $\mathcal L(K_X+S)\cong \OO _X(K_{X}+S +D)$ and
    $\mathcal L(K_X)\cong \OO _X(K_{X} +D)$.
    
    Since $(X,S)$ is plt, then $\lfloor \Delta _{X'} \rfloor \leq 0$ and so $\lceil \mu ^*D\rceil \geq \lfloor \Delta _{X'} +\mu ^*D \rfloor$.
    Thus
    \[K_{X'}+S'+\lceil \mu ^*D\rceil\geq \lfloor K_{X'}+S'+\Delta _{X'} +\mu ^*D \rfloor =\lfloor \mu ^*(K_X+S +D) \rfloor , \qquad {\rm and}\]
    \[K_{X'}+\lceil \mu ^*D\rceil\geq \lfloor K_{X'}+\Delta _{X'} +\mu ^*D \rfloor \>\lfloor \mu ^*(K_X+D) \rfloor .\]
    It follows that 
    \[\mu _* \OO _{X'}(K_{X'}+S'+\lceil \mu ^*D\rceil)=\OO _X(K_{X}+S +D), \qquad {\rm and}\]
    \[\mu _* \OO _{X'}(K_{X'}+\lceil \mu ^*D\rceil)=\OO _X(K_{X} +D).\]

    
     Consider the following exact sequence
    \[0\to \OO _{X'}(K_{X'}+\lceil \mu ^* D\rceil )\to \OO _{X'}(K_{X'}+S'+\lceil \mu ^* D\rceil )\to  \mathcal E'\to 0,\] 
    where $\mathcal E'=\OO _{X'}(K_{X'}+S'+\lceil \mu ^* D\rceil )|_{S'}$ is a line bundle.
    Since $K_{X'}+\lceil \mu ^* D\rceil\equiv _X K_{X'}+\{ -\mu ^*D \} $, where $(X',\{ -\mu ^*D \})$ is klt and $\mu$ is bimeromorphic,
    by the relative Kawamata-Viehweg vanishing theorem (see Lemma \ref{l-wdvan}), we have $R^1\mu _* \OO _{X'}(K_{X'}+\lceil \mu ^* D\rceil )=0$.
    Thus, applying $\mu _*$, we obtain an exact sequence
    \[0\to \OO _X(K_{X} +D)\to \OO _X(K_{X} +S+D)\to \mathcal E \to 0 \] where $\mathcal E =\mu _* \mathcal E'$.
   % where $\mathcal E =\mu _* \OO _{S'}(K_{S'}+\lceil \mu ^* D\rceil |_{S'})$. %Let $G=\mu _* G'$, then 
     By Lemma \ref{lem:CM-sheaves}, $\OO _X(K_{X} +D)$ and $\OO _X(K_{X} +S+D)$ are Cohen-Macaulay, and hence so is $\mathcal E$ by \cite[Corollary 2.62(3)]{Kol13}. Thus by \cite[Proposition 1.3]{Har80}, $\mathcal E$ is reflexive sheaf on $S$ of rank 1.
  
Now, since $S$ is normal, it is smooth in codimension 1, and since $\mathcal E$ is reflexive, it is locally free in codimension 1.
Then there is a big open subset $S^0\subset S$ such that $S^0$ is smooth and $\mathcal E|_{S^0}$ is locally free.
 Recall that  $\mu:X'\to X$ is given by a sequence of blow ups along smooth centers. Let $X^0\subset X$ be the open subset of $X$ obtained by removing the images of the exceptional divisors of $\mu$ whose centers have codimension $\>3$ in $X$. Thus,  over $X^0$, $\mu$ is given by a sequence of blow ups of smooth centers of codimension $2$. We may assume that $S^0=S\cap  X^0$. By classification of surface plt singularities it follows that we may also assume that there is an integer $m>0$ such that $m(K_X+S)$ and $\mathcal L ^{\otimes m}$ are Cartier on $X^0$. 
It suffices now to check that $\Delta _S=\Delta _{\rm diff}-\nu_*(\{-\mu ^* D\}|_{S'})\geq 0$ and that the natural inlcusion \[(\mathcal E|_{S^0}) ^{\otimes m}\to (\mathcal L  \otimes \OO _X(K _X+S))^{\otimes m}|_{S^0}=(\mathcal L ^{\otimes m})|_{S^0}\otimes \OO _{S_0}(m(K_S+\Delta _{\rm diff}))\] induces the equality $(\mathcal E|_{S^0}) ^{\otimes m}\cong \OO _{S^0}(m(K_S+\Delta _S))\otimes \mathcal O _X(mD)|_{S^0}$  
on a neighborhood $U_s\subset X^0$ of each point $s\in S^0$.
We may assume that $U_s$ is Stein. Cutting by general hyperplanes passing through $s$ we may assume that $\dim X=2$. The claim now follows by a standard computation for surface plt singularities as in \cite[Lemma 3.3]{HW19}.
     

\end{proof}


The following result is a generalization of Theorem \ref{t-pl} that works for generalized pairs.

\begin{theorem}\label{t-ext}
Let $(X,S+B+\bbeta )$ be a gplt pair, where $X$ is a compact K\"ahler variety such that $K_X$, $B$, and $S$ are all $\mathbb R$-Cartier. Suppose that 
\begin{enumerate}
    \item $\lrd S+B\rrd=S$ is irreducible, 
    \item there is a proper morphism $\pi:S\to T$ such that $\pi _*\OO_S=\OO_T$,
    \item $-(K_X+S+B+\bbeta_X)|_S$ is relatively K\"ahler (over $T$), and 
    \item $-S|_S$
is $\pi$-ample.
\end{enumerate}  Then there exists a bimeromorphic morphism $p:X\to Z$ such that $p|_S=\pi$ and $p|_{X\setminus S}$ is an isomorphism.
\end{theorem}
\begin{proof}
From adjunction as in Lemma \ref{lem:normality-of-plt-center} we know that $S$ is normal, $(S,B_S+\bbeta _S)$ is gklt, $B_S\geq {\rm Diff}_S(S+B)$, and  $-(K_S+B_S+\bbeta _S) $ is K\"ahler over $T$, where $(K_X+S+B+\bbeta _X)|_S=K_S+B_S+\bbeta _S$.
Since $K_X+S$ and $B$ are $\mathbb R$-Cartier, $(X,S)$ is plt. 


Choosing $\mcL=\mcO_X(-K_X-(k+1)S)$ for all $k\>1$ in \eqref{eqn:reflexive-exact} we get the following short exact sequences
 \begin{equation}\label{eqn:reflexive-exact-k}
 0\to \OO_X(-(k+1)S)\to \OO_X(-kS)\to \mathcal E_k\to 0.
 \end{equation}

Since $S$ is $\Q$-Cartier, $\OO _X(S)$ is a $\Q$-line bundle, by Proposition \ref{p-hw} it follows that $\mcE_k$ is a reflexive rank $1$ sheaf on $S$. 
Let $S_{k}$ be the thickening of $S$ defined by $\OO_{S_k}\cong \OO_X/\OO_X(-kS)$. Then we have the following short exact sequence of sheaves of rings on $S$
\[0\to \mathcal E_k \to  \OO_{S_{k+1}}\to \OO_{S_k}\to 0 .\]
We can then consider the corresponding long exact sequence of sheaves of rings on $T$ obtained by pushforward (note that identifying the topological spaces of the thickenings $|S_k|=|S|$, we may regard $\pi:|S_k|\to |T|$ as a map of topological spaces).
We claim that: 
\begin{claim}\label{c-van} 
$R^1\pi _* \mathcal E_k =0$ for $i>0$, and so $\pi_* \OO_{S_{k+1}}\to \pi_* \OO_{S_k}$ is surjective for all $k\>1$.
\end{claim}



Grant this for the time being. Since $\pi_* \OO_{S_{k+1}}\to \pi_* \OO_{S_k}$ is surjective for $k\geq 1$,  we have inclusions $T_k\hookrightarrow T_{k+1}$, where $T_k:={\rm Specan} (\pi _* \OO_{S_{k}})$. Since $\pi_* \OO_{S_1}=\pi_*\mcO_S=\mcO_T$, the $T_k$ define thickenings  
of $T$  with underlying topological space $|T|$. The natural isomorphisms $H^0(\OO _{T_k})\cong H^0(\pi _* \OO_{S_{k}})\cong H^0(\OO _{S_k})$ induce morphisms of complex spaces $\pi _k :S_k\to T_k$ (e.g. see \cite[Exercise II.2.4]{Har77}). The induced composite morphisms $S_k\to T_k\injective T_{k+1}$ factors as $S_k\to S_{k+1}\to T_{k+1}$. We would like to emphasize here that $\pi_k:S_k\to T_k$ are morphisms of non-reduced complex spaces for all $k\geq 1$, not just maps of topological spaces.
Fix $\ell\>1$ such that $A=\ell S=S_\ell$ is Cartier and let $\pi _A:A\to T_\ell$ be the induced morphism (here $\pi_A=\pi _\ell$). 


Consider the following commutative diagram 
\begin{equation}
    \xymatrixcolsep{3pc}\xymatrixrowsep{3pc}\xymatrix{
S=A_{\red}\ar[r]^{\vphi}\ar[d]_{\pi} & A\ar[d]^{\pi_A}\\
T=(T_\ell)_{\red}\ar[r]^{\psi} & T_\ell.
    }
\end{equation}
We claim that $\mcO_A(-A)$ is $\pi_A$-ample. Indeed, since $\vphi^*\mcO_A(-A)=\mcO_S(-\ell S|_S)$ is $\pi$-ample, and $\psi$ is a proper finite morphism, $\vphi^*\mcO_A(-A)$ is $\pi_A\circ\vphi$-ample. Since $\vphi:A_{\red}\to A$ is the reduction of nilpotent elements of $\mcO_A$, from an argument similar to the proof of \cite[Proposition 1.2.16(i)]{Laz04a} it follows that $\mcO_A(-A)$ is $\pi_A$-ample. 

Arguing as in the proof of Lemma \ref{l-cont}, we will show the following. 
\begin{claim}\label{c-van+}
    Replacing $A$ by a multiple, we may assume that $R^i\pi  _{A,*}\OO _{A}(-kA)=0$ for all $k\>1$.
\end{claim} 

By Theorem \ref{t-fuj74}, it then follows that there exists a proper bimeromorphic morphism $p:X\to Z$ such that $p|_{X\setminus A}$ is an isomorphism and $p|_{A}=\pi _A$.
Finally, $p|_S=(p|_A)|_S=\pi _A|_S=\pi$.


\begin{proof}[Proof of Claim \ref{c-van}]To check the claim, note that the vanishing $R^1\pi _* \mathcal E_k =0$ can be checked locally on the base $T$, and hence we may assume that $T$ is a relatively compact Stein variety. Since $\pi$ is a projective morphism,  it follows  that $\mathcal E_k=\OO _S(G_k)$ for some Weil divisor $G_k$ on $S$. 
By Proposition \ref{p-hw}, $G_k \sim _\Q -kS|_S-\Delta_k$, where $\Delta_k\>0$ is an effective $\mbQ$-divisor on $S$ such that $0\leq \Delta _k\leq {\rm Diff}_S(S+B)$.
If $S$ is $\Q$-factorial, then
\[ G_k\sim _\Q K_{ S}+B_S-\Delta _k+\bbeta _S-(K_S+B_S+\bbeta _S)-kS|_S,\]
where $\bbeta_S$ is defined by adjunction as in Subsection \ref{subsec:adjunction}, $(S, (B_S-\Delta _k)+\bbeta _S)$ is gklt by Lemma \ref{lem:normality-of-plt-center}, and $-(K_S+B_S+\bbeta _S)-kS|_S$ is relatively K\"ahler over $T$. The claim now follows from Lemma \ref{l-wdvan}.

Otherwise,  since $(S, {\rm Diff}_S(S+B))$ is  klt, there is  a small $\mbQ$-factorialization $q:\tilde S\to S$ (see \cite[Theorem 2.19]{DHY23}). Let $\tilde G_k:=q^{-1}_*G_k$, $\tilde B_S:=q^{-1}_*B_S $, $ \tilde \Delta _k:=q^{-1}_*\Delta _k$ so that $K_{\tilde S}+\tilde B_S-\tilde \Delta _k+\bbeta _{\tilde S} =q^*(K_S+B_S-\Delta _k+\bbeta _S)$
and $\tilde{\mathcal E} _k :=\OO _{\tilde S}(\tilde G_k)$. Then
\[\tilde G_k\sim _\Q K_{\tilde S}+\tilde B_S-\tilde \Delta _k+\bbeta _{\tilde S}-q^*(K_S+B_S+\bbeta _S+kS|_S),\] where 
$(\tilde  S, (\tilde B_S-\tilde \Delta _k)+\bbeta _{\tilde S})$ is  gklt and $-q^*(K_S+B_S+\bbeta _S+kS|_S)$ is relatively nef and big over $T$ (hence also over $S$).

By Lemma \ref{l-wdvan}, $R^iq_*\OO _{\tilde S}(\tilde G_k)=0$ and $R^i(\pi \circ q )_*\OO _{\tilde S}(\tilde G_k)=0$ for all $i>0$.
Since $q$ is small, then $q _*\OO _{\tilde S}(\tilde G_k)=\OO _S(G_k)$, and so by an easy spectral sequence argument we have \[R^i\pi_*\OO _S(G_k) =R^i\pi_*( q _*\OO _{\tilde S}(\tilde G_k))=0\qquad \mbox{ for all } \ i>0.\] 
\end{proof}


\begin{proof}[Proof of Claim \ref{c-van+}] 
Since the statement is local over $T_\ell$, replacing $T_\ell$ by a Stein open subset $U$ and $X$ by a neighborhood of $\pi _A^{-1}(U)$, we may assume that $T_\ell$ is Stein and the statement is equivalent to $H^i(A,\OO _A(-kA))=0$ for all $k>0$. Since $\OO _A(-A)$ is ample, by Serre vanishing, there is an integer $k_0>0$ such that $H^i(A,\OO _A(-kA))=0$ for all $k\geq k_0$ and $i>0$.
For $j\>1$, consider the short exact sequence
  \[0\to \OO _{A}(-(k+j)A)\to \OO _{(j+1)A}(-kA)\to \OO _{jA}(-kA) \to 0.\]
  Proceeding by induction we get that $H^i(\OO _{(j+1)A}(-kA))=0$ for all $k\geq k_0$, $j\>1$, and $i>0$. Replacing $A$ by $k_0A$, the claim now follows.
 \end{proof}
\end{proof}



The next result shows that flips for generalized pairs exist in all dimensions.
\begin{theorem}\label{t-flip}
Let $(X,B+\bbeta)$ be a $\mbQ$-factorial gdlt pair, where $X$ is a compact analytic variety of arbitrary dimension in Fujiki's class $\mcC$, and $f:X\to Z$ a flipping contraction. Then the flip of $f$ exists. 
\end{theorem}
\begin{proof} 
Replacing $B$ by $(1-\epsilon)B$ for some $0<\epsilon<1$, we may assume that $(X,B+\bbeta )$ is gklt. Let $\nu :X'\to X$ be a log resolution of $(X,B+\bbeta )$ such that $f\circ\nu$ is a projective morphism and $K_{X'}+B'+\bbeta_{X'}=\nu^*(K_X+B+\bbeta_X)$.
Let $\bbeta _{X'}+E=\nu ^*\bbeta _X$, then $E\geq 0$ is an effective $\mbR$-divisor by the negativity lemma.

Since $\rho (X/Z)=1$, it follows that if $K_X+B\not \equiv _Z0$, then $X\to Z$ is projective and $\bbeta _X\equiv _Z \lambda (K_X+B)$
for some $\lambda \ne 0$. For any point $z\in Z$ we fix a neighborhood $z\in W\subset Z$ such that $W$ is relatively compact, Stein and satisfies Property ({\bf P}) (see \cite[Definition 2.17 and Remark 2.18]{DHY23}). Let $X_W=f^{-1}(W)$, then we may assume that $K_{X_W}$
is a $\Q$-Cartier divisor and we let $D:=\lambda (K_{X_W}+B_W)$ where $B_W=B|_{X_W}$. We also let $X'_W=\nu ^{-1}(X_W)$, $B'_W=B'|_{X'_W}$, $E_W=E|_{X'_W}$ and $D':=\nu _W ^*D-E$. Then $D'\equiv_W \bbeta _{X'_W}$ is nef and big over $W$. Since $D'\equiv_W \bbeta _{X'_W}$ is nef and big over $W$, we may assume that $D'\equiv_W \Delta' _W\geq 0$ so that $(X'_W,B'_W+\Delta '_W)$ is sub-klt, and hence $(X_W,B_W+\Delta _W)$ is klt, where $B_W+\Delta_W:=\nu _{W,*}(B'_W+\Delta '_W)$. But then the relative log canonical model $(X^+_W,B^+_W+\Delta^+ _W)$ of $(X_W,B_W+\Delta _W)$ exists by \cite[Theorem 1.3]{DHP22} and \cite[Theorem 1.8]{Fuj22}, and it is the relative log canonical model of $(X_W,B_W+\bbeta _W)$. Since these relative log canonical models are unique, they glue together, and we obtain a relative log canonical model $(X^+,B^++\bbeta )$ of $(X,B+\bbeta )$ over $Z$. 

Suppose now that $K_X+B \equiv _Z0$. Then $K_Z+B_Z:=f_*(K_X+B)$ is klt and $Z$ has rational singularities.  Fix $z\in Z$, we will construct the flip locally over a neighborhood of $z$. We may assume that a resolution $\nu:X'\to X$ is projective over $Z$ and $X'_z$  is a divisor with simple normal crossing support. Let $f'=f\circ \nu$. Since $Z$ has rational singularities, from the short exact sequence  (\cite[\S 3, eqn. (2)]{HP16})
\[0\to \R \to \OO _{X'}\to \mathcal H _{X'}\to 0\]
it follows that $R^1f'_*\mathcal H _{X'}\to R^2f'_*\R$ is an isomorphism. Note that $ R^2f'_*\R\cong R^2f'_*\mbZ\otimes_{\mbZ}\mbR$ by the universal coefficient theorem, as $f'$ is proper. Now consider the class $[\bbeta  _{X'}]\in H^{1,1}_{\rm BC}(X')\cong H^1(
\mathcal H _{X'})$. By \cite[Lemma 12.1.1]{KM92}, there is a neighborhood $z\in W\subset Z$ and a $\R$-divisor $D'_W$ on $X'_W$ such that  the images of $\bbeta'_W:=\bbeta_{X'}|_{X'_W}$ and $D'_W$ in $R^2f'_*\R$ coincide. But then $\bbeta '_W\equiv _W D'_W$. Shrinking $W$, we may assume that it is relatively compact, Stein and satisfies Property ({\bf P}). We now argue as above and construct the flip $X\dasharrow X^+$ locally over $Z$ and then glue these together.\end{proof}

\begin{remark} 
	Even though it is not clear that the flipping and flipped contractions are projective over $Z$, the proof shows that they are locally projective over $Z$.
\end{remark}
%%%%%%%%%%%%%%%%%%%%%%%%%%%%%%%



\begin{proof}[Proof of Theorem \ref{t-pl}]
This follows immediately from Theorems \ref{t-ext} and \ref{t-flip}.
\end{proof}















%%%%%%%%%%%%%%%%%%%%%%%%%%%%%%%%%%%%%%%%%%%%%%%%%%%%%%%%%%%%%%%%%%%%%%
%%%%%%%%%%%%%%%%%%%%%%%%%%%%%%%%%%%%%%%%%%%%%%%%%%%%%%%%%%%%%%%%%%%%%%

\subsection{Existence of 3-fold flips and divisorial contractions}
Suppose that $X$ is a normal compact analytic variety in Fujiki's class $\mcC$, and $\alpha\in H^{1,1}_{\rm BC}(X)$ is nef and big but not K\"ahler, then by \cite[Theorem 2.30]{DHP22}, ${\rm Null}(\alpha)$ is non-empty. Recall that ${\rm Null}(\alpha)$ is the union of (positive dimensional) analytic subvarieties $Z\subset X$ such that $\alpha ^{\dim Z}\cdot Z=0$. Suppose that $\alpha ^\perp \cap \overline{\rm NA}(X)=R$ is an extremal ray, then we say that $R$ is divisorial if $\dim {\rm Null}(\alpha)=\dim X-1$ (i.e. ${\rm Null}(\alpha)$ contains a divisor) and $R$ is of flipping type if $\dim {\rm Null}(\alpha)<\dim X-1$ (i.e. ${\rm Null}(\alpha)$ contains no divisors). 
If $\alpha =[K_X+B+\bbeta _X+\omega]$ where $(X,B+\bbeta )$ is a glc pair and $\omega$ is K\"ahler form, then it is expected that ${\rm Null}(\alpha)$ is a closed analytic subset of $X$ {\it and} that ${\rm Null}(\alpha)$ is given a by the union of all curves $C$ such that $[C]\in R$, i.e. $\alpha\cdot C=0$. Recall that when $X$ is smooth and $\alpha $ is nef and big, then by \cite{CT15} it is known that ${\rm Null}(\alpha)= E_{nK}(\alpha)$ is a closed analytic subset of $X$. 


The result below shows that flipping contractions and flips exist in dimension $3$. 
\begin{theorem}\label{t-flipcont}  
Let $(X,B)$ be a $\Q$-factorial compact K\"ahler 3-fold klt pair, $\beta$ a K\"ahler form on $X$ such that 
\begin{enumerate}
%    \item $t\beta $ is very general in $H^{1,1}_{\rm BC}(X)$ for some $t>0$,
    \item $\alpha =[K_X+B+\beta]$ is nef and big, and
\item $\alpha ^\perp \cap \overline{\rm NA}(X)=R$ is an extremal ray of flipping type.
\end{enumerate} Then the flipping contraction $f:X\to Z$ and the flip $X\dasharrow X^+$ exist, and there is a K\"ahler class $\alpha _Z$ on $Z$ such that $\alpha=f^*\alpha _Z$.
\end{theorem}
\begin{proof} Let $\nu:X'\to X$ be a log resolution of $(X, B)$. We may assume that $\alpha ':=\nu ^*\alpha =\omega '+G'$ where $G'\geq 0$ is $\nu$-exceptional and $\omega '$ is a K\"ahler class (see the arguments in the proof of Proposition \ref{p-1}). Let $\bomega':=\overline{\omega}'$ (see \cite[Subsection 2.1]{DHY23}), and $\Delta':=\nu ^{-1}_*B+{\rm Ex}(\nu)$, then $(X',\Delta' )$ is dlt and we may write  $K_{X'}+\Delta '=\nu ^*(K_X+B)+E$ where $E\geq 0$ and $\Supp(E)=\Ex(\nu)$.
Note that $K_{X'}+\Delta '+t\omega '$ is K\"ahler for $t\gg 0$.
Thus
 the following hypothesis are satisfied:
\begin{enumerate}
\item $\alpha '$ is nef and $\omega '$ is a modified K\"ahler class, 
\item $\alpha '=\omega '+G'$, where $G'\geq 0$ is supported on $\lfloor \Delta '\rfloor$,
    \item $K_{X'}+\Delta '+t\omega '+a\alpha '$ is nef for some $t>0$ and $a>0$, and
    \item $(X',\Delta '+t\bomega ')$ is $\Q$-factorial gdlt (see \cite[Subsection 2.1]{DHY23}).
\end{enumerate}
Note that even though $\omega'$ is a K\"ahler class here, in the rest of proof we will only use the modified K\"ahler property of $\omega'$, which is preserved by steps of the MMP. 

If $S'$ is any component of $\lfloor \Delta '\rfloor$, then by adjunction we can write 
\[K_{S'}+\Delta _{S'}+t\omega' _{S'}=(K_{X'}+\Delta'+t\bomega'_{X'})|_{S'},\]
where $\Delta _{S'}:={\rm Diff}_{S'}(\Delta ')$.
Then $(S',\Delta _{S'})$ is a dlt surface and $\omega' _{S'}={\bomega'_{X'}}|_{S'}$ %$\omega' _{S'}=\overline{({\omega'}|_{S'})}$. Let $\omega_{S'}$ denote the trace of $\bomega'_{S'}$ on $S'$; then $\omega_{S'}$ 
is a big class, since $\bomega'_{X'}$ is a modified K\"ahler class. 

Let $\alpha _{S'}:=\alpha' |_{S'}$ so that $\alpha _{S'}$ is nef and define
 \[\tau _{S'}:={\rm inf}\{s\geq 0\;|\; K_{S'}+\Delta _{S'}+s\omega _{S'}+a\alpha _{S'}\ {\rm is\ nef\ for\ some\ }a>0\}.\]
 \begin{claim}\label{c-S}
     If $\tau _{S'}>0$, then $K_{S'}+\Delta _{S'}+\tau _{S'}\omega _{S'}+a'\alpha _{S'}$ is nef for $a'\gg 0$ and there is a $(K_{S'}+\Delta _{S'})$-negative extremal ray $R_{S'}$ in $\overline{\rm NA}(S')$ such that
     \[(K_{S'}+\Delta _{S'}+\tau _{S'}\omega _{S'})\cdot R_{S'}=\alpha _{S'} \cdot R_{S'}=0.\]
 \end{claim}
 \begin{proof}
     Since $\omega _{S'}$ is big and $\tau _{S'}>0$, then there are finitely many $(K_{S'}+\Delta _{S'}+\frac {\tau _{S'}}2\omega _{S'})$-negative extremal rays in $\overline{\rm NA}(S')$ (see \cite[Corollary 2.32]{DHY23}). Let $\Sigma _1,\ldots , \Sigma _k$ be curves that span those rays 
     and let \[\Lambda:=\{i\;:\; \alpha_{S'}\cdot \Sigma_i=0 \mbox{ for some } 1\<i\<k\}.\] We define
     \[\sigma _{S'}:={\rm inf}\left\{s\geq \frac{\tau _{S'}}2\;|\;(K_{S'}+\Delta _{S'}+s\omega _{S'})\cdot \Sigma _i\geq 0\mbox{ for all } i\in\Lambda\right\}.\] 
    If $\Lambda=\emptyset$, i.e. $\alpha _{S'}\cdot \Sigma _i> 0$ for $i=1,\ldots, k$, then we define $\sigma _{S'}:=\frac{\tau _{S'}}2$.
     
     Clearly 
     $\sigma _{S'}\leq \tau _{S'}$ and if $\sigma _{S'}= \tau _{S'}$, then there is a curve $\Sigma _i$ for some $i\in\Lambda$ such that $(K_{S'}+\Delta _{S'}+\tau _{S'}\omega _{S'})\cdot \Sigma _i=\alpha _{S'}\cdot \Sigma _i=0$.
     % Since $K_{S'}+\Delta _{S'}+t\omega _{S'}+a\alpha_{S'}$ is nef and $t\geq \tau _{S'}$, 
     Since $\Sigma_i$ is $(K_{S'}+\Delta_{S'}+\frac{\tau_{S'}}{2}\omega_{S'})$-negative, it follows that $(K_{S'}+\Delta_{S'})\cdot\Sigma_i<0$ and the claim holds. Therefore, it suffices to show that $\sigma _{S'}< \tau _{S'}$ is impossible. To the contrary assume that $\sigma _{S'}< \tau _{S'}$. Then for $a'\gg 0$ we have \[(K_{S'}+\Delta _{S'}+\sigma _{S'}\omega _{S'}+a'\alpha _{S'})\cdot \Sigma _i\geq 0\qquad {\rm for}\ i=1,\ldots, k.\]
     For $0<\epsilon \ll 1$ and $a'\gg 0$, we claim that $K_{S'}+\Delta _{S'}+(\tau _{S'}-\epsilon )\omega _{S'}+a'\alpha _{S'}$ is non-negative on \[\overline {\rm NA}(S')_{K_{S'}+\Delta _{S'}+\frac {\tau _{S'}}2\omega _{S'}\geq 0}.\] 
     Indeed, $\alpha _{S'}$ is nef and $K_{S'}+\Delta _{S'}+(\tau _{S'}-\epsilon )\omega _{S'}+a'\alpha _{S'}$ is a convex linear combination of $K_{S'}+\Delta _{S'}+\frac {\tau _{S'}}2\omega _{S'}$ and the nef classes $K_{S'}+\Delta _{S'}+t\omega _{S'}+a\alpha _{S'}$ and $\alpha _{S'}$. More specifically,  if $\eta _{s,a}:= K_{S'}+\Delta _{S'}+s\omega _{S'}+a\alpha _{S'}$ and $\lambda:=\frac {t-\tau _{S'}+\epsilon}{t-\tau _{S'}/2}$, then we can write \[\eta_{\tau _{S'}-\epsilon ,a'} =\lambda  \eta_{{\tau _{S'}}/2,0}+(1-\lambda) \eta_{t,a} +(a'-(1-\lambda)a)\alpha _{S'}.\]
     Note that $\lambda >0$,  $1-\lambda=\frac {\tau _{S'}/2-\epsilon}{t-\tau _{S'}/2}>0$, since $0<\eps\ll 1$, and $a'-(1-\lambda)a \geq 0$ for $a'\gg 0$.
     Since  $K_{S'}+\Delta _{S'}+(\tau _{S'}-\epsilon )\omega _{S'}+a'\alpha _{S'}$ is non-negative on the $\Sigma _i$ for all $1\<i\<k$, it is non-negative on 
     \[ \overline {\rm NA}(S')=\overline {\rm NA}(S')_{K_{S'}+\Delta _{S'}+\frac {\tau _{S'}}2\omega _{S'}\geq 0}+\sum _{i=1}^k \mathbb R^+[\Sigma _i].\]
     Thus $K_{S'}+\Delta _{S'}+(\tau _{S'}-\epsilon )\omega _{S'}+a'\alpha _{S'}$ is nef, which is a contradiction to the definition of $\tau_{S'}$, and the claim follows.
     Note that  $K_{S'}+\Delta _{S'}+\tau _{S'}\omega _{S'}+a'\alpha _{S'}$ is nef for all $a'\gg 0$, as $\alpha_{S'}$ is nef.
 \end{proof}~\\
 

Let $\tau:={\rm max}\{\tau _{S'}\}$ as $S'$ runs through all the components of  $\lfloor \Delta '\rfloor$. Then $K_{S'}+\Delta _{S'}+\tau\omega _{S'}+a'\alpha _{S'}$ is nef for each component $S'$ of  $\lfloor \Delta '\rfloor$ and $a'\gg 0$, as $K_{S'}+\Delta_{S'}+t\omega_{S'}+a'\alpha_{S'}$ is nef and $\tau_{S'}\<\tau\<t$ for every component $S'$ of $\lrd\Delta'\rrd$.

\begin{claim} If $\tau >0$, then $K_{X'}+\Delta '+\tau \omega'+a\alpha'$ is nef for some $a>0$.
%We have $\tau ={\rm inf}\{ s\geq 0|K_{X'}+\Delta '+s\omega'+a\alpha'\ {\rm is\ nef\ for\ some\ }a>0\}$.
\end{claim}
\begin{proof}
    Suppose that $K_{X'}+\Delta '+\tau \omega'+a'\alpha'$ is not nef for a given $a'\gg 0$, then by \cite[Theorem 2.36 and Remark 2.37]{DHP22} there is a subvariety $Z'\subset X'$ such that $(K_{X'}+\Delta '+\tau \omega'+a'\alpha ')|_{Z'}$ is not pseudo-effective. %Clearly $\lambda \leq \tau <t$ and so $t-\lambda>0$. 
Since
\[K_{X'}+\Delta'+\tau \omega'+a'\alpha' \equiv  K_{X'}+\Delta'+t \omega'+(a'-t+\tau)\alpha'+(t-\tau)G' \]
and $K_{X'}+\Delta '+t \omega'+(a'-t+\tau)\alpha '$ is nef for $a'\gg 0$,
 then $Z'$ is contained in the support of $G'$ and hence in a component $S'$ of $\lfloor \Delta '\rfloor$ and this contradicts what we have shown above (see Claim \ref{c-S}).
\end{proof}


If $\tau >0$, then from Claim \ref{c-S} it follows that there is a component $S'$ of $\lfloor \Delta '\rfloor $ and a $(K_{S'}+\Delta _{S'})$-negative  extremal ray $R_i=\mbR^+\cdot[\Sigma_i]$ such that  $(K_{X'}+\Delta '+\tau \omega')\cdot R_i=\alpha'\cdot R_i=0$. Let $\pi : S'\to T$ denote the corresponding contraction. Since $\omega'\cdot R_i>0$ and  $\alpha '\cdot R_i=(\omega '+G ')\cdot R_i=0$, it follows that $G'\cdot \Sigma_i<0$. Thus $\Sigma_i$ is contained in a component of $G'$, say $S''$ such that $S''\cdot\Sigma_i<0$. Since $\Supp(G')=\lrd\Delta'\rrd$, $\Sigma_i$ is contained in a $(K_{S''}+\Delta_{S''})$-negative extremal face of $\NA(S'')$, where $K_{S''}+\Delta_{S''}:=(K_{X'}+\Delta')|_{S''}$. This face contains a negative extremal ray $R''_i=\mbR^+\cdot[\Sigma''_i]$ such that $S''\cdot R''_i<0$. Thus replacing $S'$ by $S''$ and $R_i$ by $R''_i$ we may assume that $S'\cdot R_i<0$, and hence $-S'|_{S'}$ is $\pi$-ample.

Now write $\Delta'=S'+\Delta''$ so that $S'\not\subset \Supp(\Delta'')$. Since $(X', \Delta'+\bomega')$ is gdlt, for $0<\eps\ll 1$ we have $(X', S'+(1-\eps)\Delta''+\bomega')$ is gplt and $(K_{X'}+S'+(1-\eps)\Delta''+\omega')\cdot R_i<0$. Thus by Theorem \ref{t-ext}, there exists 
a bimeromorphic morphism $p : X' \to Z'$ such that $p|_{S'} = \pi$ and $p|_{X'\setminus S'}$ is an isomorphism.
\begin{claim}\label{clm:contraction}
   The morphism $p$ is either a divisorial or flipping contraction, and in the latter case the corresponding flip exists. 
\end{claim}
\begin{proof}
    Since $-(K_{X'}+\Delta ')$ is $p$-ample, $Z'$ has rational singularities by \cite[Lemma 2.41]{DH20}. Let $\gamma \in H^{1,1}_{\rm BC}(X')$ such that $\gamma \cdot R_i=0$, then $\gamma \cdot C=0$ for any $p$-exceptional curve. This is because all $p$-exceptional curves are contained in $S'$ and $\rho (S'/T)=1$.
    By \cite[Lemma 3.3]{HP16}, there is a class $\gamma _{Z'}\in H^{1,1}_{\rm BC}(Z')$ such that $p^*\gamma _{Z'}=\gamma$. It follows that $\rho (X'/Z')=1$.
    This shows that 
 $p$ is either a divisorial or flipping contraction. By Theorem \ref{t-flip}, if it is a flipping contraction, then the corresponding flip exists. 
 \end{proof}

We may then replace $X'$ by the image of this flip or divisorial contraction. Observe that, at this stage $X'$ is in Fujiki's class $\mcC$ and not necessarily K\"ahler, however, by Lemma \ref{lem:fujiki-to-kahler-surface} the components of $\lrd\Delta'\rrd$ are still K\"ahler surfaces with $\mbQ$-factorial rational singularities. So we can repeat the above arguments. 


\begin{claim}\label{c-m} Fix $0<\epsilon  \ll 1$. After finitely many steps, we obtain a bimeromorphic map $\phi: X'\dasharrow X^+$ such that for any $0<t\leq \epsilon$ \[K_{X^+}+\Delta ^++t\omega ^+ +a\alpha ^+=\phi _*(K_{X'}+\Delta'+t\omega'+a\alpha ')\] is nef for some $a>0$ (depending on $t$).
\end{claim} 
\begin{proof}
    Properties (1-4) stated at the beginning of our proof continue to be satisfied after each flip or divisorial contraction. So we may repeat the procedure obtaining an $\alpha '$-trivial, $(K_{X'}+\Delta ')$-minimal model program. Since 
    termination of flips hold by Theorem \ref{thm:termination}, after finitely many steps we may assume that $\tau=0$. Thus we obtain an $\alpha '$-trivial $(K_{X'}+\Delta ')$-MMP $X'\dasharrow X^+$ such that $K_{X^+}+\Delta ^++\epsilon \omega ^++a\alpha ^+$ is nef for some $\epsilon >0$, $a>0$ 
    and \[{\rm inf}\{ s\geq 0\;|\; K_{X^+}+\Delta ^++s\omega^++a\alpha^+\ {\rm is\ nef\ for\ some\ }a>0\}=0.\]
   The claim now follows easily by taking convex linear combinations. 
\end{proof}


Let $\psi:X\dasharrow X^+$ and $\phi:X'\dasharrow X^+$ be the induced bimeromorphic maps,  $U:=X\setminus\Null(\alpha)$ and $X'_{U}:=\nu^{-1}(U)$. Next we will show that every step of this MMP is vertical over $U$.




\begin{claim}\label{c-p}
The map $\phi _U=:\phi |_{X'_U}$ is a proper bimeromorphic map defined over $U$, i.e. there exist proper bimeromorphic morphisms $X'_U\to U$ and $X^+_U\to U$ commuting with $\phi_U$.
\end{claim} 
\begin{proof}
  
Suppose that $X'=:X^0\dasharrow X^1\dasharrow \ldots \dasharrow X^n=:X^+$ is the MMP obtained above and $\phi ^i:X'\dasharrow X^i$ are the induced bimeromorphic maps. Proceeding by induction, it suffices to show that if
$\phi ^i_U:=\phi ^i|_{X'_U}:X'_U\dasharrow X^i_U$ is a proper bimeromorphic map defined over $U$, then so is $\phi ^{i+1}_U$. We will denote by $\nu ^i_U:X^i_U\to U$ the corresponding proper morphism.
Since $X^i\dasharrow X^{i+1}$ is an $\alpha ^i$-trivial flip or divisorial contraction,  if $f^i: X^i\to Z^i$ is the corresponding contraction, and $C$ is a contracted curve, then $\alpha ^i\cdot C=0$. Suppose that $C_U:=C\cap X^i_U\ne \emptyset$. 
Let $p:W\to X'$, $q:W\to X^i$ be the normalization of the graph of $\phi ^i$, then \[0=C\cdot \alpha ^i=q^{-1}_*C\cdot q^* \alpha ^i=p_*q^{-1}_*C\cdot  \alpha '=\nu _* p_*q^{-1}_*C\cdot \alpha .\] If $C_U$ is not $\nu ^i_U$ vertical, then $\nu _* p_*q^{-1}_*C\ne 0$ and so $\nu _* p_*q^{-1}_*C\cdot \alpha \ne 0$ which is a contradiction. Thus every such $C_U$  is $\nu ^i_U$ vertical.  It follows that $f^i_U:=f^i|_{X^i_U}$ is a (proper bimeromorphic) morphism over $U$ and hence $X^{i+1}_U\to U$ is a proper bimeromorphic map over $U$. 
  
\end{proof}

   \begin{claim}\label{c-s}
       The divisors contracted by $\phi: X'\dasharrow X^+$ coincide with the set of $\nu: X'\to X$ exceptional divisors and hence $\psi:X\dasharrow X^+$ is a small bimeromorphic map of $\Q$-factorial varieties.
   \end{claim}
   \begin{proof} 
   We will first show that ${\rm Supp} N(K_{X'}+\Delta '+a\alpha ')={\rm Ex}(\nu)$ for any $a\gg 0$. Recall that for any pseudo-effective $(1,1)$ class $\gamma$, $N(\gamma)$ is defined as the negative part of the Boucksom-Zariski decomposition of $\gamma$, see Definition \ref{def:mod-nef}.\\
   Now since $\alpha$ is big, $K_X+B+a\alpha$ is big for $a\gg 0$. Moreover,  since $K_{X'}+\Delta '+a\alpha ' =\nu ^*(K_X+B+a\alpha)+E$, where ${\rm Supp}(E)={\rm Ex}(\nu)$, from \cite[Lemma A.5]{DHY23} it follows that ${\rm Supp} N(K_{X'}+\Delta '+a\alpha ' )\supset {\rm Ex}(\nu)$.
Since $\alpha '=G'+\omega '$ where $\omega '$ is K\"ahler and $\Supp(G')={\rm Ex}(\nu)$,  $\omega ^*:=\omega '+\frac 1 a (K_{X'}+\Delta ')$ is also K\"ahler for $a\gg 0$ and so $K_{X'}+\Delta '+a\alpha ' \equiv a\omega ^*+aG'$. Thus ${\rm Supp} N(K_{X'}+\Delta '+a\alpha ')\subset {\rm Supp} (N(G'))\subset {\rm Ex}(\nu)$ and hence ${\rm Supp} N(K_{X'}+\Delta '+a\alpha ')={\rm Ex}(\nu)$ for any $a\gg 0$.
Since $\omega'$ is a modified K\"ahler class, ${\rm Supp} N(K_{X'}+\Delta '+a\alpha '+\epsilon \omega ')={\rm Ex}(\nu)$ for all $0<\epsilon \ll 1$ (e.g. see the proof of \cite[Claim 3.18]{DHY23}). 
     Since $\phi: X'\dasharrow X^+$ is  a $(K_{X'}+\Delta '+\epsilon \omega '+a\alpha ')$-MMP, then $K_{X^+}+\Delta ^++\epsilon \omega ^+ +a\alpha ^+$ is nef, and by \cite[Theorem A.11]{DHY23}, $\phi$ contracts ${\rm Ex}(\nu)$.
   \end{proof}
By \cite[Lemma 2.27]{DH20}, $\psi_*:H^{1,1}_{\BC}(X)\to H^{1,1}_{\BC}(X^+)$ is an isomorphism. Let $R^\perp\subset N^1(X)$ be the codimension 1 subspace of classes $\gamma \in N^1(X)$ such that $\gamma \cdot R=0$. Let  $W$ be the normalization of the graph of $\psi:X\bir X^+$ and $p:W\to X$ and $q:W\to X^+$ the induced morphisms. Since $\psi$  is not a morphism, then there is a curve $\tilde C\subset W$ such that $p_*\tilde C=0$ and $C^+=q_* \tilde C\ne 0$. It follows that $\alpha ^+\cdot C^+=0$. Let $R^+\subset N_1(X^+)$ be the ray spanned by $C^+$.
If $\gamma \in R^\perp$, then $\nu ^* \gamma \cdot \Sigma \equiv 0$ for any curve $\Sigma$ such that $\alpha ' \cdot \Sigma =0$ where $\alpha '=\nu ^*\alpha $.
Since the minimal model program $X'\dasharrow X^+$ is $\alpha'$-trivial, i.e. it only contracts $\alpha'$-trivial curves, it follows that it is also $\gamma ':=\nu ^* \gamma $-trivial and hence that $\gamma ^+\cdot C^+=0$ where $\gamma ^+=\phi _* \gamma '\in H^{1,1}_{\rm BC}(X^+)$.
Therefore, $\psi _*(R^\perp)=(R^+)^\perp$, and hence every $p$-exceptional curve $\tilde C'$ satisfies $q_*\tilde C'\in R^+$.
Since $\psi_*$ is an isomorphism and $(K_X+B)\cdot R\ne 0$, then $(K_{X^+}+B^+)\cdot R^+\ne 0$. Suppose that $(K_{X^+}+B^+)\cdot R^+< 0$, then 
$(K_{X^+}+B^++t\omega^+)\cdot R^+< 0$ for $0<t\ll 1$. Since 
$\alpha ^+\cdot R^+=0$, $(K_{X^+}+B^++t\omega^++a\alpha ^+)\cdot R^+< 0$ for any $a>0$, contradicting Claim \ref{c-m}. Therefore $(K_{X^+}+B^+)\cdot R^+>0$.

Let $\eta :=\alpha +\delta (K_{X}+B)$ for some $0<\delta \ll 1$ and  $\eta ^+:=\psi _* \eta\in H^{1,1}_{\BC}(X^+)$.
Consider $E:=p^*\eta -q^*\eta ^+$, then we claim that $-E$ is $p$-ample. Indeed, if  $\tilde C\subset W$ $p$-vertical curve, then $q_*\tilde C\ne 0$, and as we have seen above that $[q_*\tilde C]\in R^+$, so $\eta ^+\cdot q_*\tilde C>0$ and hence $-E\cdot \tilde C>0$. It follows by the negativity lemma that $E\geq 0$ and the support of $E$ contains $\Ex(p)$. But then ${\rm Supp}(E)=\Ex(p)$. Clearly $\Supp(E)\subset \Ex(q)$, as $\psi$ is an isomorphism in codimension $1$. 
We claim that $\Ex(q)=\Supp(E)=\Ex(p)$. Indeed, let $\tilde C\subset W$ be a $q$-vertical curve. Then $p_*C\neq 0$ and in fact $p_*C\subset \Null(\alpha)$ as $\psi:X\bir X^+$ is an isomorphism on $X\setminus\Null(\alpha)$. Thus $E\cdot \tilde C=\eta\cdot p_*\tilde C=(\alpha+\delta(K_X+B))\cdot p_*\tilde C<0$, and hence $\Ex(q)\subset \Supp(E)$. 

We will now show that $-E|_{{\rm Supp}(E)}$ is ample. Let $E_i$ be a component of $\Ex(p)=\Ex(q)$. We have $C_i:=p(E_i)$ and $C_i^+:=q(E_i)$ are curves and $\theta:=(p|_{E_i}, q|_{E_i}):E_i\to C_i\times C_i^+$ is finite. Let $\tilde C_i$ (resp. $\tilde C_i^+$) be a curve on $E_i$ dominating $C_i\times \{x\}$ (resp. $\{x\}\times C_i^+$) for general $x\in C_i^+$ (resp. $x\in C_i$), then $q_* \tilde C_i=0$ (resp. $p_*\tilde C_i^+=0$). Thus  $[C_i^+]\in R^+$ and $C_i\cdot \alpha =0$ so that  $[C_i]\in R$. It follows that $-\eta|_{C_i^+}$ (resp. $\eta|_{C_i}) $ is K\"ahler. Since $\theta $ is finite, then  $-E|_{E_i}\equiv (q^*\eta ^+-p^*\eta)|_{E_i}=\theta^*(-\eta|_{C_i},\eta^+|_{C^+_i})$ is ample and so 
$-E|_{{\rm Supp}(E)}$ is ample. Perturbing the coefficients of $E$ slightly we may assume that $E$ is a $\mbQ$-Cartier divisor. But then we can contract $E$ to a point by Lemma \ref{l-cont}. 

Let $g:W\to Z$ be the induced bimeromorphic morphism, then we claim that we have induced morphisms $f:X\to Z$ and $f^+:X^+\to Z$. To this end, if $C\subset W$ is a $q$ (respectively $p$) exceptional curve, then $C$ is contained in ${\rm Supp}(E)=\Ex(p)=\Ex(q)$. %and so $\eta \cdot p_*C<0$ (respectively,  $\eta ^+\cdot q_* C>0$).Since at least one of $p_*C,q_*C$ is non-zero, then \[ E\cdot C=(p^*\eta -q^*\eta ^+)\cdot  C <0\] and so $C$ is contained in the support of $E$. 
Since $g$ contracts $E$ to a point, then
$g_*C=0$ and so by the rigidity lemma $X\to Z$ and $X^+\to Z$ are morphisms.

Note that by construction $f:X\to Z$ contracts a non-empty set of $\alpha$-trivial curves (i.e. curves in $R$). We claim that in fact $f$ contracts every curve in $R$. Let $C\subset X$ be a curve in $R$ so that $\alpha \cdot C=0$. If $C$ is not contracted by $f$, then $p$ is an isomorphism over the general points of $C$ and we let $\tilde C:=p^{-1}_*C$ and $C^+:=q_* \tilde C$. Then $\alpha ^+\cdot C^+=0$.
Since $K_{X^+}+B^++t\omega ^++a\alpha ^+$ is nef for any $0<t\ll 1$ and some $a\gg 0$, then
\[(K_{X^+}+B^++t\omega ^+)\cdot C^+=(K_{X^+}+B^++t\omega ^++a\alpha ^+)\cdot C^+\geq 0\]
and taking the limit we have $(K_{X^+}+B^+)\cdot C^+\geq 0$.
It follows that
\[0\leq E\cdot \tilde C=(p^*\eta -q^*\eta ^+)\cdot \tilde C=(K_X+B)\cdot C-(K_{X^+}+B^+)\cdot C^+<0.\]
This is impossible and so $C$ is contracted by $f$.
 It follows that  $f$ is $(K_X+B)$-negative. By \cite[Lemma 2.41]{DH20}, $Z$ has rational singularities and by \cite[Lemma 3.3]{HP16} $f^*H^{1,1}_{\rm BC}(Z)=R^\perp\subset H^{1,1}_{\rm BC}(X)$.
Thus $\rho (X/Z)=1$.
Finally, since $f$ contracts the set of $\alpha$-trivial curves, we have $\alpha =f^*\alpha _Z$. 

We will now check that $\alpha _Z$ is a K\"ahler class. 
To this end, let $V\subset Z$ be a positive dimensional subvariety of $Z$.    Note that $ \dim f(\Ex(f))=0$. Let $V'$ be the strict transform of $V$. Then $(\alpha_{Z})^{\dim V}\cdot V={\alpha}^{\dim V}\cdot V'>0$ and $V'$ is not contained in ${\rm Null}(\alpha)$.  Thus from \cite[Lemma 6.3]{DH20} or more generally from \cite[Theorem 2.30]{DHP22} it follows that $\alpha_{Z}$ is a K\"ahler class. 
Since $\psi _*$ is an isomorphism, then $\rho (X^+/Z)=1$.
We have already seen that $f^+$ contracts a $(K_{X^+}+B^+)$-positive curve and so $K_{X^+}+B^+$ is $f^+$-ample.
Finally, it is  easy to see that $(f^+)^*\alpha _Z+\delta (K_{X^+}+\Delta ^+)$ is K\"ahler for $0<\delta\ll 1$, and hence $X^+$ is a K\"ahler $3$-fold.
   

\end{proof}~\\


Next we prove the existence of divisorial contractions in dimension $3$.
\begin{theorem}\label{t-divcont} Let $(X,B)$ be a $\Q$-factorial compact klt 3-fold pair, $\beta$ a K\"ahler form on $X$ such that 
\begin{enumerate}
%    \item $t\beta $ is very general in $H^{1,1}_{\rm BC}(X)$ for some $t>0$,
    \item $\alpha =[K_X+B+\beta]$ is nef and big, and
\item $\alpha ^\perp \cap \overline{\rm NA}(X)=R$ is an extremal ray of divisorial type.
\end{enumerate} Then the divisorial contraction $f:X\to Z$ exists and there is a K\"ahler class $\alpha _Z$ on $Z$ such that $\alpha \equiv f^*\alpha _Z$.
\end{theorem}
\begin{proof} Recall that by definition $\alpha ^\perp \cap \overline{\rm NA}(X)=R$ is an extremal ray of divisorial type if and only if $\dim {\rm Null}(\alpha)=2$. 
We begin by observing the following.
\begin{claim}
    There is an irreducible divisor $S$ such that ${\rm Null}(\alpha)=S$, $S$ is Moishezon, $S$ is covered by $\alpha$-trivial curves, and $S\cdot R<0$.
\end{claim}
\begin{proof}
Let $\omega$ be a K\"ahler class on $X$. Since $\alpha$ is big, $\alpha-\eps\omega$ is big for $0<\epsilon \ll 1$, and we have the Boucksom-Zariski decomposition $\alpha -\epsilon \omega =\sum s_jS_j+P$ where $P$ is modified nef and big and in particular the restriction of $P$ to any surface is pseudo-effective. Let $S\subset\Null(\alpha)$ be a divisor (a surface) and $S'\to S$ the minimal resolution. Then $\alpha |_{S'}$ is nef but not big (as $(\alpha |_{S'})^2=0$), therefore $S$ coincides with some component of $\sum s_jS_j$, which for simplicity we denote by $S_1$.
    We claim that $S$ is Moishezon. To this end, consider $b={\rm mult}_{S}(B)$ and \[\left(1+\frac {1-b}{s_1}\right)\alpha \vert_{S'}=\left(K_X+B+\beta +\left(\frac {1-b}{s_1}\right)\left(\sum s_jS_j+P+\epsilon \omega\right)\right)\Big\vert_{S'}.\]
     Since ${\rm mult}_{S}\left(B+\frac {1-b}{s_1}\sum s_jS_j\right)=1$, we have $\left(K_X+B +\frac {1-b}{s_1}\sum s_jS_j\right)\big\vert_{S'}=K_{S'}+B_{S'}$ where $B_{S'}\geq 0$.
       %Suppose first that $\alpha |_{S'}\equiv 0$. 
       Then 
    \begin{equation}\label{eqn:restriction-to-resolution}
         \left(1+\frac {1-b}{s_1}\right)\alpha |_{S'}=K_{S'}+B_{S'}+\left(\beta +\frac {1-b}{s_1}(P+\epsilon \omega)\right)\Big\vert_{S'}.
    \end{equation} 
    Since $B_{S'}\geq 0$, $P|_{S'}$ is pseudo-effective, $(\beta +\frac {1-b}{s_1}(P+\epsilon \omega))|_{S'}$ is big and $\alpha |_{S'}$ is not big, it follows that $K_{S'}$ is not pseudo-effective, and hence $S'$ is projective by classification and $H^2(S', \mcO_{S'})=0$ (e.g. see \cite[Lemma 2.40]{DH20}). In particular, $S$ is Moishezon. 
    
    We claim that $S$ is covered by an analytic family of $\alpha$-trivial curves $\{C_t\}$.
    Fix $A_{S'}$ an ample divisor on $S'$. Since $H^2(S', \mcO_{S'})=0$, $\alpha'|_{S'}$ is represented by an $\mbR$-divisor. 
    Since $\alpha |_{S'}$ is nef, it is in the cone $\overline {NM}_1(S')$ of movable curves. By \cite[Theorem 1.3]{Ara10} and \cite[Theorem 1.9]{Das20}, for any $\epsilon >0$, we have
a decomposition \[\alpha |_{S'}\equiv C_{\epsilon}+\sum _{i=1}^k\lambda _iM_i,\]
where $C_{\epsilon}\in (\overline{NE}_1(S')_{(K_{S'})\geq 0}+\overline {NM}_1(S'))_{(K_{S'}+\epsilon A_{S'})\geq 0}$, $\lambda _i>0$, $(K_{S'}+\epsilon A_{S'})\cdot C_{\epsilon}\geq 0$ and the $M_i$ are movable curves. If $\alpha |_{S'}\cdot M_i=0$, then the pushforward of $M_i$ onto $S$ defines a movable curve on $S$ which is $\alpha$-trivial and the claim follows. Therefore, we may assume that $\alpha |_{S'}\cdot M_i>0$ for all $i$.
Now, $0=(\alpha |_{S'})^2=\alpha |_{S'}\cdot (C_{\epsilon}+\sum _{i=1}^k\lambda _iM_i)$ and since $\alpha |_{S'}$ is nef then $k=0$, i.e. $\alpha |_{S'}\equiv C_{\epsilon}$. But then $\alpha |_{S'}\cdot (K_{S'}+\epsilon A_{S'})\geq 0$. Taking the limit as $\epsilon \to 0$, it follows that
$\alpha |_{S'}\cdot K_{S'}\geq 0$. Since $\alpha |_{S'}$ is nef, then $\alpha |_{S'}\cdot \gamma\geq 0$ for any pseudo-effective class $\gamma$. Thus $\alpha |_{S'}\cdot B_{S'}\geq 0$, $\alpha |_{S'}\cdot P|_{S'}\geq 0$, and $\alpha |_{S'}\cdot \omega |_{S'}\geq 0$. It follows from \eqref{eqn:restriction-to-resolution} and $(\alpha'|_{S'})^2=0$ that all these terms are equal to $0$. But then $\alpha |_{S^\nu}\cdot \omega |_{S^\nu}=0$ where $S^\nu\to S$ is the normalization. Since $\omega |_{S^\nu}$ is K\"ahler, then $\alpha |_{S^\nu}\equiv 0$. This contradicts the inequality above $\alpha |_{S'}\cdot M_i>0$, and so the claim holds.
   
    
Let $\{C_t\}_{t\in T}$ be an analytic family of $\alpha$-trivial curves on $S$ covering $S$ as proved in the claim above. Then $[C_t]\in R$ for all $t\in T$ and we have
\[0>-\epsilon \omega\cdot C_t =(\alpha -\epsilon \omega)\cdot C_t =\left(\sum s_jS_j+P\right)\cdot C_t\geq s_{1}S\cdot C_t\]
as $S_i\cdot C_t\geq 0$ for $i\neq 1 $ and $S=S_1$. Thus $S\cdot C_t<0$, and hence $S\cdot R<0$ as $R=\mbR^+\cdot[C_t]$. In particular, if $C$ is any curve with $[C]\in R$, then $S\cdot C<0$ and hence $S={\rm Null}(\alpha)$.
\end{proof}
 Let $b:={\rm mult}_S(B)$ and $\Delta:=B+(1-b)S$. 
 \begin{claim} There exists a dlt model $\mu :X'\to X$ of $(X,\Delta)$. In particular 
\begin{enumerate}
    \item $(X',\Delta ')$ is $\Q$-factorial and dlt, where $\Delta':=\mu ^{-1}_*\Delta +{\rm Ex}(\mu)$, 
    \item $K_{X'}+\Delta '$ is nef over $X$ so that $\mu ^*(K_X+\Delta )-(K_{X'}+\Delta ')\geq 0$, and
    \item $K_{X'}+\Delta '=\mu ^*(K_X+B)+(1-b)S'+\sum a_jE_j$, where $S'=\mu ^{-1}_*S$ and $a_j>0$ for all $j$. 
    \end{enumerate}
    \end{claim}
    \begin{proof}
        Let $\mu :X'\to X$ be a projective log resolution and $\Delta':=\mu ^{-1}_*\Delta +{\rm Ex}(\mu)$, then $(X',\Delta ')$ is dlt. We may run the $(K_{X'}+\Delta ')$-MMP over $X$ (the existence of this MMP is well known since $\dim X=3$ and $\mu$ is projective, but see also \cite[Theorem 1.4]{DHP22} for the case $\dim X>3$). By \cite[Theorem 3.3]{DO23}, we may assume that the corresponding MMP terminates and replacing $(X',\Delta ')$ by the output of this MMP, we may assume that $K_{X'}+\Delta '$ is nef over $X$.
        By the negativity lemma, $\mu ^*(K_X+\Delta )-(K_{X'}+\Delta ')\geq 0$ and so (1) and (2) hold.
        (3) follows since $(X,B)$ is klt.
        
    \end{proof}

Let $\eta'$ be a K\"ahler class on $X'$ and note that by \cite[Lemma 2.27]{DH20}, there exists a $\mu$-exceptional $\mbR$-divisor $F$ such that $\eta'+F\equiv _X0$. Thus $-F\equiv_X \eta '$ is ample over $X$ and by the negativity lemma $F\geq 0$. In particular, ${\rm Ex}(\mu)={\rm Supp}(F)$. Perturbing the coefficients of $F$ we may assume that $F$ is a $\mbQ$-divisor. It follows that the non-klt locus of $(X',\Delta ')$ is equal to $S'\cup {\rm Ex}(\mu)$. By (2), the image of ${\rm nklt}(X',\Delta ')$ is contained in ${\rm nklt}(X,\Delta )=S$. Thus $X'_U\to U$  is an isomorphism where $U:=X\setminus S=X\setminus\Null(\alpha)$.

We will now run the $\alpha'$-trivial $(K_{X'}+\Delta')$-minimal model program, where $\alpha '=\mu ^*\alpha$.

If $C$ is a curve not contained in $W':=X'\setminus X'_U$, then $\alpha'\cdot C>0$. It follows easily that every $\alpha' $-trivial step of the $(K_{X'}+\Delta ')$-MMP will involve only curves contained in the complement of $X'_U$, and hence the restriction of this MMP to $X'_U$ is an isomorphism. 

Let $n(\alpha )$ be the nef dimension of the restriction of $\alpha$ to the normalization $S^\nu\to S$, see \cite[Definition 2.7]{BCE02}. Since $S$ is a Moishezon surface covered by a family of $\alpha$-trivial curves (showed above), it follows that $0\leq n(\alpha )\leq 1$. If $n(\alpha )=0$, then $\alpha |_{S^\nu}\equiv 0$ and if $n(\alpha )=1$, then the nef reduction map $S^\nu \to T$ is a morphism to a smooth (projective) curve (see \cite[2.4.4]{BCE02}). Observe that if $n(\alpha )=0$, then the nef reduction map $S^\nu \to T$ is a morphism to a point and in this case the existence of the divisorial contraction $f:X\to Z $ follows from Corollary \ref{c-1}. We however give an alternate argument below which proves both of the cases: $n(\alpha)=0$ and $n(\alpha)=1$,  simultaneously. 


 If $\gamma\in H^{1,1}_{\BC}(X)$, then we will say that $\gamma\equiv_\alpha 0$ if and only if $\gamma\cdot C=0$ for every $\alpha$-trivial curve $C\subset X$, i.e. $\gamma\cdot C=0$ whenever $\alpha \cdot C=0$ for a curve $C\subset X$.
We  have that $(K_X+B-sS)\cdot R=0$ for some $s>0$ (as both $K_X+B$ and $S$ are $R$-negative), and hence $K_X+B-sS\equiv_\alpha 0$.
It follows that 
$\mu ^* (K_X+B-sS)\equiv _{\alpha'}0$ and hence that \[K_{X'}+\Delta'\equiv_{\alpha'} \Theta':=(1-b)S'+\sum a_jE_j+s\mu ^*S\geq 0,\] where %$\mathcal E'=(1-b)S'+\sum c_jE_j+s\mu ^*S\geq 0$ and 
${\rm Supp}(\Theta')=\lfloor \Delta '\rfloor$. 

\begin{claim}
    We can run the $\alpha '$-trivial $(K_{X'}+\Delta ')$-MMP 
    \[X'=X^0\dasharrow X^1 \dasharrow \ldots \dasharrow X^m\]
    so that if $\phi^i:X'\dasharrow X^i$, then
\begin{enumerate}
\item each flip and divisorial contraction is $\alpha'$-trivial, hence $\alpha ^i=\phi ^i_*\alpha'$ is nef,
\item  $\phi ^i|_{X'_U}$ is an isomorphism over $U$ and we denote by $X^i_U$ its image,
\item $(X^i,\Delta ^i:=\phi ^i_*\Delta ')$ is dlt, and $\lfloor \Delta ^i\rfloor \subset W^i:=X^i\setminus X^i_U$, 
\item $K_{X^i}+\Delta ^i-\Theta^i\equiv_{\alpha ^i} 0$ and ${\rm Supp}(\Theta^i)=\lfloor \Delta ^i\rfloor$, and
    \item if $S^m$ is a component of $\lfloor \Delta ^m\rfloor$ then every $K_{S^m}+\Delta _{S^m}:=(K_{X^m}+\Delta ^m)|_{S^m}$ negative extremal ray $R^m$ is $\alpha _{S^m}:=\alpha _{X^m}|_{S^m}$-positive, i.e. $\alpha _{S^m}\cdot R^m>0$.
\end{enumerate}
    
\end{claim}
\begin{proof}
We will proceed by induction on $i$. Assume that we have already constructed $X'\dasharrow X^1 \dasharrow \ldots \dasharrow X^i$ with the required properties.
Suppose that there is a component $S^i$ of $\lfloor \Delta ^i\rfloor$ and a $K_{S^i}+\Delta _{S^i}:=(K_{X^i}+\Delta ^i)|_{S^i}$-negative extremal ray $R^i=\R^+[\Sigma ^i]$ which is $\alpha _{S^i}:=\alpha _{X^i}|_{S^i}$ non-positive and hence trivial: $\alpha _{S^i}\cdot R_i=0$.
Since $K_{X^i}+\Delta ^i\equiv_{\alpha ^i} \Theta ^i$ and ${\rm Supp}(\Theta^i)=\lfloor \Delta ^i\rfloor$, then $\Theta^i \cdot \Sigma ^i<0$ and hence, there is a component $\bar S^i$ of $\lfloor \Delta ^i\rfloor$ such that $\bar S^i\cdot \Sigma ^i<0$ 
and so $\Sigma ^i\subset \bar S^i$. Clearly  $\Sigma ^i$ is also $K_{\bar S^i}+\Delta _{\bar S^i}:=(K_{X^i}+\Delta ^i)|_{\bar S^i}$-negative and hence is contained in an $\alpha _{\bar S^i}$-trivial $(K_{\bar S^i}+\Delta _{\bar S^i})$-negative extremal face. Let $\bar \Sigma ^i$ be a curve spanning an extremal ray in this face, we may assume that $\bar S^i\cdot \bar \Sigma ^i<0$. Let $\pi :\bar S^i\to V$ be the corresponding extremal contraction,
then $-\bar S^i $ and $-(K_{\bar S^i}+\Delta _{\bar S^i})$ are relatively ample over $V$ and so by Theorem \ref{t-ext} there is a contraction $p:X^i\to Z^i$ such that $p|_{\bar S^i}=\pi$ and $p|_{X^i\setminus \bar S^i}$ an isomorphism. If $p$ is a flipping contraction, then the flip exists by Theorem \ref{t-flip}.
We let $X^i\dasharrow X^{i+1}$ be the induced bimeromorphism (be it a flip or divisorial contraction).
By construction $p$ is $\alpha ^i$-trivial and $p|_{X^i_U}$ is the identity. In particular, since $Z^i$ has rational singularities (this follows from \cite[Lemma 2.41]{DH20}), $\alpha ^i\equiv p^*\alpha _{Z^i}$ where $\alpha _{Z^i}$ is nef.  (1-3)$_{i+1}$ are easily seen to hold.
Since $K_{X^i}+\Delta ^i-\Theta^i\equiv_{\alpha ^i} 0$ and $p$ is $\alpha ^i$-trivial, it also follows that \[K_{Z^i}+\Delta _{Z^i}-\Theta_{Z^i}:=p_*(K_{X^i}+\Delta ^i-\Theta^i)\equiv_{\alpha _{Z^i}} 0\] and hence, pulling back to $X^{i+1}$ we have $K_{X^{i+1}}+\Delta ^{i+1}-\Theta^{i+1}\equiv_{\alpha ^{i+1}} 0$ and (4)$_{i+1}$  holds.
By termination of flips \cite[Theorem 3.3]{DO23}, after finitely many steps, we obtain $X\dasharrow X^m$ such that 
(5) holds.
\end{proof}
Now recall the nef reduction morphism $S^\nu\to T$.  Note that if $P$ is a component of $\mu ^* S$, then there is an induced morphism $\sigma _P:P\to T$ such that for any curve $C$ on $P$ we have $\alpha '\cdot C=0$ if and only if $C$ is vertical over $T$ i.e. $\sigma _{P,*}C=0$. Note that if $n(\alpha|_{S^\nu})=0$, then $T$ is a point and every curve is vertical over $T$. Since $\phi ^i:X'\dasharrow X^i$ is $\alpha '$-trivial, it follows easily that if $\phi ^i_*P=P^i\ne 0$, then the induced bimeromorphic map $P\dasharrow P^i$ is defined over $T$ and in particular $P^i\dasharrow T$ is a morphism.
\begin{claim}
    The map $\phi^m:X'\dasharrow X^m$ contracts every component of $\mu ^*S$ dominating $T$, and hence if $n(\alpha )=0$, then $\phi :X\dasharrow X^m$ is a bimeromorphic map that contracts $S$ (and no other divisors) and if $n(\alpha )=1$, then $\phi :X\dasharrow X^m$ contracts only $S$ and may extract some divisors of $X'$ which are contained in $\mu ^{-1}(S)$ that do not dominate $T$.
\end{claim}
\begin{proof}
    Recall that $F'\geq 0$ is a $\mu$-exceptional $\Q$-divisor such that $-F'$ is $\mu$-ample. Since $-S|_{S^\nu}\equiv_\alpha  \frac {-1}s(K_X+B)|_{S^\nu}$ %$\equiv_\alpha \frac {1}s\omega _{S^\nu}$
    is ample over $T$, then $\Xi':=\mu ^*S+\epsilon F'$ is a $\Q$-divisor such that $-\Xi'|_{S'}$ and $-\Xi'|_{E_j}$ are ample over $T$ for every component $E_j$ of ${\rm Ex}(\mu)$ and all $0<\epsilon \ll 1$. 
    
    We claim that $\phi^m$ is $\Xi'$ non-positive, i.e. if $p:W\to X'$ and $q:W\to X^m$ is the normalization of the graph of $\phi^m$, then we claim that $p^*\Xi'-q^*(\phi ^m_*\Xi')\geq 0$. Recall that $p$ and $q$ are isomorphisms over $X_{U}'=\mu^{-1}U=\mu^{-1}(X\setminus S)$, since so is $\phi^m$. We now check that $-p^* \Xi'$ is nef over $X^m$. To this end consider a curve $C$ on $W$ such that $q_*C=0$, then $p_*C\ne 0$ is contained in $\Supp(\mu^*S)$ and $\alpha '\cdot p_*C=\alpha ^m\cdot q_*C=0$.
    But then $p_*C$ is vertical over $T$ and so $-p^* \Xi'\cdot C=-\Xi'\cdot p_*C>0$; in particular, $-p^*\Xi'$ is nef over $X^m$.
    %As $-\mathcal F'$ is ample on every component of $\mu ^{-1}(S)=X'\setminus X'_U$, it follows that $-p^* \mathcal F'$ is nef over $X^m$. 
    Let $\Xi^m:=q_* p^*\Xi'=\phi^m_*\Xi'$, then $q^*\Xi^m-p^*\Xi'$ is $q$-exceptional and $q$-nef so that by the negativity lemma, $p^*\Xi'-q^*\Xi^m\geq 0$; this proves our claim.\\
    
    Suppose now that $\phi^m _*(\mu ^{*}S)$ dominates $T$ i.e. that $\phi^m _*(\mu ^{*}S)\to T$ is surjective and let $\lambda$ be the smallest constant such that $(\Theta^m-\lambda \Xi^m)_\textsubscript{hor}\leq 0$,  where $(\ldots )_\textsubscript{hor}$ denotes the horizontal components over $T$ i.e. those components that dominate $T$. Then there is a component $E^m$ of $\phi^m _*(\mu ^{*}S)$ dominating $T$  such that ${\rm mult }_{E^m}(\Theta^m-\lambda \Xi^m)=0$, and  there is a family  $C^m_t$ (contained in the fibers over $T$) of $\alpha ^m$-trivial curves covering $E^m$. We claim that $(K_{X^m}+\Delta ^m)\cdot C^m_t\geq 0$.  
    If this is not the case, then \[ [C^m_t]\in (\alpha ^m) ^\perp\cap \overline{\rm NA}(E^m)_{K_{E^m}+\Delta _{E^m}<0 }\] where $K_{E^m}+\Delta _{E^m}=(K_{X^m}+\Delta ^m)|_{E^m}$. Since $\alpha^m$ is nef, by the cone theorem on $E^m$, there is a $(K_{E^m}+\Delta _{E^m})$-negative $\alpha^m$-trivial extremal ray, contradicting (5) above.
    If $C_t$ is the strict transform on $E=(\phi^m) ^{-1}_*E^m$, then 
    \[0\geq (\Theta^m-\lambda \Xi^m)\cdot C^m_t= (K_{X^m}+\Delta ^m-\lambda \Xi^m) \cdot C^m_t\geq -\lambda \Xi^m \cdot C^m_t\geq -\lambda \Xi'\cdot C_t>0\]
    where the first inequality follows by our definition of $\lambda$, the second as $(K_{X^m}+\Delta ^m-\Theta ^m)\equiv _{\alpha ^m}0$, the third was observed above, the fourth as $p^*\Xi'- q^*\Xi^m\geq 0$ and ${\rm mult}_{E^m}(p^*\Xi'- q^*\Xi^m)=0$ and the last as $-\Xi'|_E$ is ample over $T$. This is the required contradiction and the claim is proven.
\end{proof}
We will now show that in fact $\phi ^m$ contracts $\mu ^*S$. To this end, it suffices to show that $\Theta^m=0$ or equivalently that $q^*\Theta^m=0$.
Since $q^*\Theta^m$ is effective and exceptional over $X$, by the negativity lemma it suffices to show that $q^*\Theta^m$ is nef over $X$. Let $C\subset W$ be a curve such that $\mu _*p_*C=0$. If $q_*C=0$, then $C\cdot q^*\Theta^m=0$. If instead $q_*C\ne 0$, then
       $\alpha ^m\cdot q_*C=\alpha \cdot \mu _* p_*C=0$. Clearly $q^*\Theta^m\cdot C\geq 0$ if $q_*C$ is not contained in any component $P$ of the support $\Theta^m$ (and hence of the support of $\lfloor \Delta ^m\rfloor$). So assume that $q_*C$ is contained in a component $P$ of $\Theta^m$. 
       Since $\Theta^m\equiv _{\alpha ^m}K_{X^m}+\Delta ^m$ and $(K_{X^m}+\Delta ^m)|_P$ is non-negative on $\alpha ^m$-trivial curves (by (5) above), then $q^*\Theta^m\cdot C=(K_{X^m}+\Delta ^m)\cdot q_*C \geq 0$. Therefore $q^*\Theta^m$ is nef over $X$ and hence $\Theta^m=0$ as observed above.

We will now show that the induced bimeromorphic map $\phi: X\dasharrow X^m$ is a morphism. Let $p:W\to X$ and $q:W\to X^m$ be a resolution of the graph of $\phi$. Pick $\omega ^m$ a K\"ahler class on $X^m$ and $E_1,\ldots , E_r$ the $p$-exceptional divisors generating $H^{1,1}_{\BC}(W)/p^*H^{1,1}_{\BC}(X)$ as in \cite[Lemma 2.27]{DH20}. Then we may write $q^*\omega ^m+\sum e_iE_i\equiv _X0$ for some $e_i\in\mbR$, and hence there is a $\omega\in H^{1,1}_{\BC}(X)$ such that $p^*\omega \equiv q^*\omega ^m+\sum e_iE_i$. %Note that $\sum e_iE_i\geq 0$ by the negativity lemma.
Pick $r\in \R$ such that $(\omega +rS)\cdot R=0$, then $p^*(\omega +rS)\equiv q^*\omega ^m+\sum e_iE_i+rp^*S$. Let $C$ be a $q$-exceptional curve. If $p_*C=0$, then $(q^*\omega ^m+\sum e_iE_i+rp^*S)\cdot C=(\omega+rS)\cdot p_*C=0$. If $p_*C\neq 0$, then 
$p^*\alpha \cdot C=q^*\alpha ^m\cdot C=0$ and so $p_*C\in R$ and hence $(\omega +rS)\cdot p_*C=0$ so that $(q^*\omega ^m+\sum e_iE_i+rp^*S)\cdot C=0$. Recall that $\Theta^m=0$, in particular, $\phi$ contracts $S$ and doesn't extract any divisor. Thus
we have shown that $\sum e_iE_i+rp^*S$ is  $q$-exceptional and $q$-numerically trivial. Then by the negativity lemma $\sum e_iE_i+rp^*S=0$ and so $p^*(\omega +rS)\equiv q^*\omega ^m$. But then, if $C$ is $p$-exceptional and not $q$-exceptional, we have
\[0=C\cdot p^*(\omega +rS)=C\cdot q^*\omega ^m=q_*C \cdot \omega ^m>0\]
which is impossible. Since no such curves exist, then $\phi: X\dasharrow X^m$ is a morphism which we will now denote by $f:X\to Z$.

As we have seen above, $f|_U$ is an isomorphism and $f$ contracts a non-empty set of $\alpha$-trivial curves. Since $\overline{\rm NA}(X)\cap \alpha ^\perp=R$, then $f$ contracts the set of all curves in $R$ (and no other curves).
Since $K_X+B+\beta \equiv _Z 0$, then
$-(K_X+B)$ is $f$-ample. Then from \cite[Lemma 2.41]{DH20} it follows that $Z$ has rational singularities.
By \cite[Lemma 3.3]{HP16}, $\rho (X/Z)=1$ and $\alpha \equiv f^* \alpha _Z$ for some $\alpha _Z\in H^{1,1}_{\rm BC}(Z)$. By \cite[Theorem 2.30]{DHP22}, $\alpha _Z$ is K\"ahler
if and only if for any positive dimensional subvariety  $W\subset Z$,  $\alpha _Z^{\dim W}\cdot W>0$ holds.
If $W\not \subset f(S)$, then the strict transform $W'\subset X$ is not contained in $S$ and so $\alpha _Z^{\dim W}\cdot W=\alpha ^{\dim W'}\cdot W'>0$. Therefore, we may assume that $W=f(S)$ is a (projective) curve. Since $f$ is a projective morphism, $S$ is projective, and thus there is a (projective curve) $W'\subset S$ such that $f(W')=W$. Then $\alpha\cdot W'>0$, since by our construction above a curve in $X$ is contracted by $f$ if and only it is $\alpha$-trivial. Thus by the projection formula $\alpha_Z\cdot W>0$; this completes our proof.

\end{proof}


\begin{proof}[Proof of Theorem \ref{thm:contractions}]
This follows from Theorems \ref{t-flipcont} and \ref{t-divcont}.
\end{proof}

\section{Cone Theorem and Minimal Models}\label{sec:cone-and-mm}
In this section we will establish the cone theorem and existence of minimal models for klt pairs $(X, B)$ of dimension $3$ such that $K_X+B$ is pseudo-effective.\\ 


Let $K\subset \R ^n$ be a compact convex set. An \textit{extremal point} of $K$ is a point  $x$ such that if $x=tx_1+(1-t)x_2$ for $x_i\in K$, then $x=x_1$ or $x=x_2$; equivalently $K\setminus \{x\}$ is convex. An \textit{exposed point} of $K$ is a point $x\in K$ such that $\{x\}=K\cap \{f=0\}$ where $\{f=0\}$ is a supporting hyperplane for $K$. Every exposed point is extremal but not vice versa. We let ${\rm ext}(K)$ be the set of extremal points, ${\rm exp}(K)$ the set of exposed points and we denote by ${\rm conv}(X)$ the convex hull of a set $X\subset \R ^n$ and ${\rm cl}(X)$ the closure of $X$. For a detailed a discussion on these topics see \cite[Chapter 1]{HW20}. We need the following result:
\begin{theorem}\label{t-conv}    Let $K\subset \R ^n$ be a non-empty compact convex set. Then
\[{\rm cl(conv(exp}(K)))={\rm conv( ext}(K))=K.\]  
In particular, the sets ${\rm exp(K)}$ and ${\rm ext(K)}$ are both non-empty.
\end{theorem}

\begin{proof}
    The equality ${\rm conv( ext}(K))=K$ is well known and due to Minkowski (see \cite[Theorem 1.21]{HW20}), and the equality ${\rm cl(conv(exp}(K)))=K$ follows from \cite[Theorem 1.23]{HW20}.
    \end{proof}
    
   \begin{proof}[Proof of Theorem \ref{t-cone}] Let $K$ be a compact convex slice of $\overline{\rm NA}(X)$, i.e. the cone generated by $K$ is equal to $\NA(X)$, and $R$ an extremal ray corresponding to an exposed point $\{x\}\in K$. Then by definition, there is a $(1,1)$ class $\alpha\in H^{1,1}_{\BC}(X)$ such that $K$ is contained in the half-space $\alpha_{\>0}:=\{\gamma\in N_1(X): \alpha\cdot\gamma\>0\}$ and $\{x\}=K\cap\alpha^\bot$. In particular, $R=\overline{\rm NA}(X)\cap \alpha ^\perp$. Since $\alpha$ is non-negative on $\overline{\rm NA}(X)$, it is nef. If $(K_X+B)\cdot R<0$, then for $0<a\ll 1$, it follows easily that $\beta :=\alpha -a(K_X+B)$ is positive on $\overline{\rm NA}(X)\setminus\{0\}$, and hence K\"ahler (e.g. see the proof of \cite[Claim 3.25]{DHY23}). Thus $\alpha =a(K_X+B+\frac 1 a \beta)$. Replacing $\alpha$ by $\frac 1 a\alpha$ and letting $\omega = \frac 1 a \beta$ we may assume that $\alpha \equiv K_X+B+\omega$, where $\omega$ is a K\"ahler class. By Theorems \ref{t-flipcont} and \ref{t-divcont}, there is a corresponding flipping or divisorial contraction and hence $R=\R ^+[\Gamma ]$ for some curve $\Gamma$ on $X$. By \cite[Theorem 4.2]{DO23}, we may assume that $\Gamma$ is a rational curve and $0>(K_X+B)\cdot \Gamma \geq -6$. Since $X$ is a K\"ahler variety, by a Douady space argument there are at most countably many such extremal rays $\{\mbR^+\cdot[\Gamma _i]\}_{i\in I}$.
 Let $V:=\overline{\rm NA}(X)_{K_X+B \geq 0}+\sum _{i\in I}\mathbb R ^+[\Gamma _i]$. Following \cite[Lemma 6.1]{HP16} it suffices to show that $\overline{\rm NA}(X) = \overline V$. To the contrary assume that $\NA(X)\setminus \overline V\neq\emptyset$. Since $\overline V$ is a convex set (as $V$ is so) and $\NA(X)$ is generated by $K$, from Theorem \ref{t-conv} it follows that there is an exposed point $x\in K$ not contained in $\overline V$. Let $R$ be the $(K_X+B)$-negative extremal ray corresponding to the exposed point $x\in K$. Then from what we have seen above, there is a rational curve $\Gamma$ such that $0<-(K_X+B)\cdot \Gamma\<6$ and $R=\mbR^+\cdot[\Gamma]$, but all such extremal rays are contained in $\overline V$ by definition of $V$, this is contradiction.
 
 \end{proof}

A corollary of the cone and contraction theorems is the following
\begin{corollary}\label{cor:dlt-contraction}
    Let $(X, B)$ be a $\mbQ$-factorial compact K\"ahler $3$-fold dlt pair such that $K_X+B$ pseudo-effective. If $R$ is a $(K_X+B)$-negative extremal ray, then there is a projective bimeromorphic morphsim $f:X\to Z$ with connected fibers such that a curve $C\subset X$ is contracted by $f$ if and only if $[C]\in R$. 
\end{corollary}

\begin{proof}
    By the cone Theorem \ref{t-cone}, there is a K\"ahler class $\omega$ such that $\alpha=K_X+B+\omega$  is nef and big and $\alpha^\bot\cap\NA(X)=R$. Choose $0<\eps\ll 1$ so that $\omega+\eps B$ is K\"ahler. Thus replacing $B$ by $(1-\eps)B$ and $\omega$ by $\omega+\eps B$ we may assume that $(X, B)$ is a klt pair, and hence the contraction follows from Theorem \ref{thm:contractions}.
\end{proof}






\begin{proof}[Proof of Theorem \ref{t-mmp-mfs}] 
Suppose that $K_X+B$ is not nef, then by Theorem \ref{t-cone} and its proof, there is a $(K_X+B)$-negative extremal ray $R$ with nef supporting class $\alpha$ such that $R=\overline{\rm NA}(X)\cap \alpha ^\perp$, and the corresponding flipping or divisorial contractions exist by Theorems \ref{t-flipcont} and \ref{t-divcont}. Let $(X,B)\dasharrow (X_1,B_1)$ be the corresponding flip or divisorial contraction. We may replace $(X,B)$ by $(X_1,B_1)$ and repeat the above procedure. By termination of flips (see \cite[Theorem 3.3]{DO23}), after finitely many steps we obtain the required minimal model. 
\end{proof}


\bibliographystyle{hep}
\bibliography{4foldreferences}
\end{document}


\newcommand{\etalchar}[1]{$^{#1}$}
\begin{thebibliography}{BHPVdV04}

\bibitem[Ara10]{Ara10}
C.~Araujo, \textsl{ The cone of pseudo-effective divisors of log varieties
  after {B}atyrev},
\newblock Math. Z. \textbf{ 264}(1), 179--193 (2010).

\bibitem[AT84]{AT84}
V.~Ancona and V.~V. Tan, \textsl{ On the blowing down problem in {${\bf
  C}$}-analytic geometry},
\newblock J. Reine Angew. Math. \textbf{ 350}, 178--182 (1984).

\bibitem[BCE{\etalchar{+}}02]{BCE02}
T.~Bauer, F.~Campana, T.~Eckl, S.~Kebekus, T.~Peternell, S.~a. Rams,
  T.~Szemberg and L.~Wotzlaw,
\newblock A reduction map for nef line bundles,
\newblock in \textsl{ Complex geometry ({G}\"{o}ttingen, 2000)}, pages 27--36,
  Springer, Berlin, 2002.

\bibitem[BCHM10]{BCHM10}
C.~Birkar, P.~Cascini, C.~D. Hacon and J.~McKernan, \textsl{ Existence of
  minimal models for varieties of log general type},
\newblock J. Amer. Math. Soc. \textbf{ 23}(2), 405--468 (2010).

\bibitem[BHPVdV04]{BHPV04}
W.~P. Barth, K.~Hulek, C.~A.~M. Peters and A.~Van~de Ven,
\newblock \textsl{ Compact complex surfaces}, volume~4 of \textsl{ Ergebnisse
  der Mathematik und ihrer Grenzgebiete. 3. Folge. A Series of Modern Surveys
  in Mathematics [Results in Mathematics and Related Areas. 3rd Series. A
  Series of Modern Surveys in Mathematics]},
\newblock Springer-Verlag, Berlin, second edition, 2004.

\bibitem[BK23]{BK23}
F.~Bernasconi and J.~Koll\'{a}r, \textsl{ Vanishing {T}heorems for
  {T}hree-folds in {C}haracteristic {$p>5$}},
\newblock Int. Math. Res. Not. IMRN (4), 2846--2866 (2023).

\bibitem[{Bou}02]{Bou02ar}
S.~{Boucksom}, \textsl{ {On the volume of a line bundle}},
\newblock arXiv Mathematics e-prints , math/0201031 (January 2002),
  {math/0201031}.

\bibitem[Bou04]{Bou04}
S.~Boucksom, \textsl{ Divisorial {Z}ariski decompositions on compact complex
  manifolds},
\newblock Ann. Sci. \'Ecole Norm. Sup. (4) \textbf{ 37}(1), 45--76 (2004).

\bibitem[Bru06]{Bru06}
M.~Brunella, \textsl{ A positivity property for foliations on compact
  {K}\"{a}hler manifolds},
\newblock Internat. J. Math. \textbf{ 17}(1), 35--43 (2006).

\bibitem[BS95]{BS95}
M.~C. Beltrametti and A.~J. Sommese,
\newblock \textsl{ The adjunction theory of complex projective varieties},
  volume~16 of \textsl{ De Gruyter Expositions in Mathematics},
\newblock Walter de Gruyter \& Co., Berlin, 1995.

\bibitem[CHP16]{CHP16}
F.~Campana, A.~H\"oring and T.~Peternell, \textsl{ Abundance for {K}\"ahler
  threefolds},
\newblock Ann. Sci. \'Ec. Norm. Sup\'er. (4) \textbf{ 49}(4), 971--1025 (2016).

\bibitem[CP97]{CP97}
F.~Campana and T.~Peternell, \textsl{ Towards a {M}ori theory on compact
  {K}\"{a}hler threefolds. {I}},
\newblock Math. Nachr. \textbf{ 187}, 29--59 (1997).

\bibitem[CT15]{CT15}
T.~C. Collins and V.~Tosatti, \textsl{ K\"{a}hler currents and null loci},
\newblock Invent. Math. \textbf{ 202}(3), 1167--1198 (2015).

\bibitem[Das20]{Das20}
O.~Das, \textsl{ Finiteness of log minimal models and nef curves on 3-folds in
  characteristic {$p>5$}},
\newblock Nagoya Math. J. \textbf{ 239}, 76--109 (2020).

\bibitem[Das21]{Das21}
O.~Das, \textsl{ Kawamata-{V}iehweg vanishing theorem for del {P}ezzo surfaces
  over imperfect fields of characteristic {$p > 3$}},
\newblock Osaka J. Math. \textbf{ 58}(2), 477--486 (2021).

\bibitem[DH20]{DH20}
O.~{Das} and C.~{Hacon}, \textsl{ {The log minimal model program for K{\"a}hler
  $3$-folds}},
\newblock arXiv e-prints , arXiv:2009.05924 (September 2020), {2009.05924}.

\bibitem[DHP22]{DHP22}
O.~{Das}, C.~{Hacon} and M.~{P{\u{a}}un}, \textsl{ {On the $4$-dimensional
  minimal model program for K{\"a}hler varieties}},
\newblock arXiv e-prints , arXiv:2205.12205 (May 2022), {2205.12205}.

\bibitem[DHY23]{DHY23}
O.~{Das}, C.~{Hacon} and J.~I. {Y{\'a}{\~n}ez}, \textsl{ {MMP for Generalized
  Pairs on K{\"a}hler 3-folds}},
\newblock arXiv e-prints , arXiv:2305.00524 (April 2023), {2305.00524}.

\bibitem[DO23]{DO23}
O.~{Das} and W.~{Ou}, \textsl{ {On the Log Abundance for Compact K{\"a}hler
  threefolds}},
\newblock Manuscripta Math.  (January 2023).

\bibitem[Fuj07]{Fuj07}
O.~Fujino,
\newblock Special termination and reduction to pl flips,
\newblock in \textsl{ Flips for 3-folds and 4-folds}, volume~35 of \textsl{
  Oxford Lecture Ser. Math. Appl.}, pages 63--75, Oxford Univ. Press, Oxford,
  2007.

\bibitem[Fuj21]{Fuj21}
O.~Fujino, \textsl{ Minimal model theory for log surfaces in {F}ujiki's class
  {$\mathcal{C}$}},
\newblock Nagoya Math. J. \textbf{ 244}, 256--282 (2021).

\bibitem[{Fuj}22]{Fuj22}
O.~{Fujino}, \textsl{ {Minimal model program for projective morphisms between
  complex analytic spaces}},
\newblock arXiv e-prints , arXiv:2201.11315 (January 2022), {2201.11315}.

\bibitem[Fuj75]{Fuj74}
A.~Fujiki, \textsl{ On the blowing down of analytic spaces},
\newblock Publ. Res. Inst. Math. Sci. \textbf{ 10}, 473--507 (1974/75).

\bibitem[Gra62]{Gra62}
H.~Grauert, \textsl{ \"{U}ber {M}odifikationen und exzeptionelle analytische
  {M}engen},
\newblock Math. Ann. \textbf{ 146}, 331--368 (1962).

\bibitem[Har77]{Har77}
R.~Hartshorne,
\newblock \textsl{ Algebraic geometry},
\newblock Graduate Texts in Mathematics, No. 52, Springer-Verlag, New
  York-Heidelberg, 1977.

\bibitem[Har80]{Har80}
R.~Hartshorne, \textsl{ Stable reflexive sheaves},
\newblock Math. Ann. \textbf{ 254}(2), 121--176 (1980).

\bibitem[HP15]{HP15}
A.~H\"oring and T.~Peternell, \textsl{ Mori fibre spaces for {K}\"ahler
  threefolds},
\newblock J. Math. Sci. Univ. Tokyo \textbf{ 22}(1), 219--246 (2015).

\bibitem[HP16]{HP16}
A.~H\"oring and T.~Peternell, \textsl{ Minimal models for {K}\"ahler
  threefolds},
\newblock Invent. Math. \textbf{ 203}(1), 217--264 (2016).

\bibitem[HW19]{HW19}
C.~Hacon and J.~Witaszek, \textsl{ On the rationality of {K}awamata log
  terminal singularities in positive characteristic},
\newblock Algebr. Geom. \textbf{ 6}(5), 516--529 (2019).

\bibitem[HW20]{HW20}
D.~Hug and W.~Weil,
\newblock \textsl{ Lectures on convex geometry}, volume 286 of \textsl{
  Graduate Texts in Mathematics},
\newblock Springer, Cham, [2020] \copyright 2020.

\bibitem[Kaw92]{Kaw92}
Y.~Kawamata, \textsl{ Termination of log flips for algebraic {$3$}-folds},
\newblock Internat. J. Math. \textbf{ 3}(5), 653--659 (1992).

\bibitem[KM92]{KM92}
J.~Koll{\'a}r and S.~Mori, \textsl{ Classification of Three-Dimensional Flips},
\newblock Journal of the American Mathematical Society \textbf{ 5}, 533--704
  (1992).

\bibitem[KM98]{KM98}
J.~Koll{\'a}r and S.~Mori,
\newblock \textsl{ Birational geometry of algebraic varieties}, volume 134 of
  \textsl{ Cambridge Tracts in Mathematics},
\newblock Cambridge University Press, Cambridge, 1998,
\newblock With the collaboration of C. H. Clemens and A. Corti, Translated from
  the 1998 Japanese original.

\bibitem[Kol13]{Kol13}
J.~Koll\'{a}r,
\newblock \textsl{ Singularities of the minimal model program}, volume 200 of
  \textsl{ Cambridge Tracts in Mathematics},
\newblock Cambridge University Press, Cambridge, 2013,
\newblock With a collaboration of S\'{a}ndor Kov\'{a}cs.

\bibitem[Laz04]{Laz04a}
R.~Lazarsfeld,
\newblock \textsl{ Positivity in algebraic geometry. {II}}, volume~49 of
  \textsl{ Ergebnisse der Mathematik und ihrer Grenzgebiete. 3. Folge. A Series
  of Modern Surveys in Mathematics [Results in Mathematics and Related Areas.
  3rd Series. A Series of Modern Surveys in Mathematics]},
\newblock Springer-Verlag, Berlin, 2004,
\newblock Positivity for vector bundles, and multiplier ideals.

\bibitem[LM22]{LM22}
S.~Lyu and T.~Murayama, \textsl{ The relative minimal model program for
  excellent algebraic spaces and analytic spaces in equal characteristic zero},
\newblock arXiv preprint arXiv:2209.08732  (2022).

\bibitem[Nak87]{Nak87}
N.~Nakayama,
\newblock The lower semicontinuity of the plurigenera of complex varieties,
\newblock in \textsl{ Algebraic geometry, {S}endai, 1985}, volume~10 of
  \textsl{ Adv. Stud. Pure Math.}, pages 551--590, North-Holland, Amsterdam,
  1987.

\bibitem[Pet98]{Pet98}
T.~Peternell, \textsl{ Towards a {M}ori theory on compact {K}\"{a}hler
  threefolds. {II}},
\newblock Math. Ann. \textbf{ 311}(4), 729--764 (1998).

\bibitem[Pet01]{Pet01}
T.~Peternell, \textsl{ Towards a {M}ori theory on compact {K}\"{a}hler
  threefolds. {III}},
\newblock Bull. Soc. Math. France \textbf{ 129}(3), 339--356 (2001).

\bibitem[Siu74]{Siu74}
Y.~T. Siu, \textsl{ Analyticity of sets associated to {L}elong numbers and the
  extension of closed positive currents},
\newblock Invent. Math. \textbf{ 27}, 53--156 (1974).

\end{thebibliography}
